\chapter{First-Order Partial Derivatives and the Gradient}

A function of one variable has a graph in the plane and a derivative that is
the slope of a tangent line. A function of two variables has a graph that is
a \emph{surface} in space, and there is no single slope any more: standing
at a point of a hillside, the ground rises steeply in one direction, falls
in another, and is level in a third. The first task of multivariable
calculus is to make that observation quantitative.

The device is to look one direction at a time. Freeze $y$ at a value and let
only $x$ move: the surface is cut by a vertical plane, the cut is an
ordinary curve, and the slope of \emph{that} curve is the partial derivative
$\partial f/\partial x$. This is the \emph{ceteris paribus} principle of
economics, and it is why partial derivatives are the natural language for
marginal products and marginal utilities.

Collecting the partial derivatives into one vector gives the gradient. It
turns out to carry all first-order information at once: it points in the
direction of steepest ascent, its length is the steepest slope available,
and it stands perpendicular to the level curve through the point. Those
three facts are proved below rather than asserted --- the third one in
particular is what the Lagrange method rests on several chapters later.

\section{Core Rules}

\begin{rulebox}{Topological vocabulary of domains}
\begin{itemize}[leftmargin=1.6em,itemsep=0.3ex]
  \item The \emph{open ball} (or \emph{neighbourhood}) of radius $r>0$ around $\mathbf{x}_0\in\R^n$ is
        $B_r(\mathbf{x}_0)=\{\mathbf{x}\in\R^n :
        \|\mathbf{x}-\mathbf{x}_0\|<r\}$.
  \item $\mathbf{x}_0\in M$ is an \emph{interior point} of $M$ if some open ball around
        $\mathbf{x}_0$ lies entirely in $M$.
  \item $M$ is \emph{open} if \emph{every} point of $M$ is
        an interior point.
  \item $\mathbf{x}_0$ is a \emph{boundary point} of
        $M$ if every ball around $\mathbf{x}_0$ contains points of $M$
        \emph{and} points of its complement.
  \item $M$ is \emph{closed} if it contains all
        its boundary points; equivalently, the complement
        $\R^n\setminus M$ is open.
  \item $M$ is \emph{bounded} if
        $M\subseteq B_r(\mathbf{0})$ for some radius $r$.
  \item $M\subseteq\R^n$ is \emph{compact} exactly when
        $M$ is closed \emph{and} bounded.
\end{itemize}
Economic examples: the open positive orthant $\{(x,y): x>0,\ y>0\}$
(the natural Cobb--Douglas domain) is open and unbounded; a budget set
$\{\mathbf{x}\in\R^n : \mathbf{p}\cdot\mathbf{x}\le m,\
\mathbf{x}\ge\mathbf{0}\}$ is closed and bounded, hence compact.
\end{rulebox}

\begin{rulebox}{Why it matters}
Differentiability (and hence every partial derivative below) is
defined at \emph{interior} points of the domain, where $f$ can be
probed in all directions. And by the \emph{Weierstrass theorem}, a
continuous function on a compact set attains its maximum and minimum:
this is why budget-constrained optimisation problems are
guaranteed to have solutions.
\end{rulebox}

\begin{rulebox}{Definition}
Let $f:\R^n\to\R$. The partial derivative of $f$ with respect to $x_i$ at
$\mathbf{x}=(x_1,\ldots,x_n)$ is
\[
  \pd{f}{x_i}(\mathbf{x}) \;=\; \lim_{h\to0}
    \frac{f(x_1,\ldots,x_i+h,\ldots,x_n)-f(x_1,\ldots,x_n)}{h} .
\]
Only the $i$-th slot changes; all other coordinates sit still. It is the
ordinary one-variable derivative of the function
$t\mapsto f(x_1,\ldots,t,\ldots,x_n)$, evaluated at $t=x_i$.
\end{rulebox}

\begin{rulebox}{What it measures: a slice of the surface}
Fix $y=y_{0}$. The vertical plane $y=y_{0}$ cuts the surface $z=f(x,y)$
along a curve, and along that curve only $x$ varies, so the curve is the
graph of the one-variable function $g(x)=f(x,y_{0})$. Then
\[
  \pd{f}{x}(x_{0},y_{0})=g'(x_{0}),
\]
the ordinary slope of the cut curve at $x_{0}$. Symmetrically,
$\partial f/\partial y$ is the slope of the cut made by the plane
$x=x_{0}$. A partial derivative is therefore never a slope \emph{of the
surface}; it is the slope of the surface \emph{along one chosen direction},
and different directions give different numbers. What the surface does in
all the remaining directions is the subject of the gradient below and of the
directional derivative.
\end{rulebox}

\begin{rulebox}{Practical computation}
To compute $\partial f/\partial x_i$, differentiate $f$ with respect to
$x_i$ and treat every other variable as a constant. All single-variable
rules (power, product, quotient, chain) apply unchanged. Worked example, for
$f(x,y)=x^{3}y^{2}+\mathrm{e}^{xy}-\ln y$:
\[
  \pd{f}{x}=3x^{2}y^{2}+y\,\mathrm{e}^{xy},
  \qquad
  \pd{f}{y}=2x^{3}y+x\,\mathrm{e}^{xy}-\frac{1}{y} .
\]
In the first computation $y$ is a number: $y^{2}$ is a constant factor, and
the chain rule on $\mathrm{e}^{xy}$ produces the inner derivative $y$. In
the second, $x$ is the constant and $\ln y$ finally contributes.
\end{rulebox}

\begin{rulebox}{Notation}
The following all refer to the same object:
\[
  \pd{f}{x} = f_x = \partial_x f = D_x f .
\]
\end{rulebox}

\begin{rulebox}{Gradient}
The vector of all first-order partial derivatives is
\[
  \grad f(\mathbf{x}) \;=\;
  \begin{pmatrix} f_{x_1}(\mathbf{x}) \\ f_{x_2}(\mathbf{x}) \\ \vdots \\
    f_{x_n}(\mathbf{x}) \end{pmatrix} .
\]
\end{rulebox}

\begin{rulebox}{Gradient and derivative: column and row}
The gradient is a \textbf{column} vector. The \textbf{derivative} of
$f\colon\R^{n}\to\R$ at $\mathbf{x}$, written $Df(\mathbf{x})$, is the same
list of partials arranged as a $1\times n$ \textbf{row}:
\[
  Df(\mathbf{x})=\bigl(f_{x_1}(\mathbf{x})\ \ f_{x_2}(\mathbf{x})\ \ \cdots\ \
  f_{x_n}(\mathbf{x})\bigr)=\grad f(\mathbf{x})^{T} .
\]
Same numbers, different shape. The row is the form that multiplies a column
of changes $\mathbf{h}$:
$Df(\mathbf{x})\,\mathbf{h}=\grad f(\mathbf{x})\cdot\mathbf{h}
=\sum_{i}f_{x_i}h_{i}$, which is the total differential of Chapter~12 and
the linear part of the tangent plane below. In Chapter~13 these rows are
stacked into the Jacobian matrix.
\end{rulebox}

\begin{rulebox}{Functions of $n$ variables written with sums}
Many economic functions are sums over goods, periods or inputs. To
differentiate with respect to one fixed variable $x_{k}$, ask which terms of
the sum contain $x_{k}$; every other term is a constant and drops out.
\begin{itemize}[leftmargin=1.6em,itemsep=0.3ex]
  \item \textbf{One function of each variable.} For
        $f=\sum_{i=1}^{n}g(x_{i})$ only the term $i=k$ survives:
        \[
          \pd{f}{x_{k}}=g'(x_{k}) .
        \]
        So $\sum_{i}x_{i}^{3}$ has $\partial f/\partial x_{k}=3x_{k}^{2}$.
  \item \textbf{The index inside the term.} $\sum_{i=1}^{n}x_{i}^{\,i}
        =x_{1}+x_{2}^{2}+\cdots+x_{n}^{n}$: the term containing $x_{k}$ is
        $x_{k}^{k}$, so $\partial f/\partial x_{k}=k\,x_{k}^{k-1}$.
        Likewise $\sum_{j=1}^{n}j/x_{j}^{4}=\sum_{j}j\,x_{j}^{-4}$ gives
        $\partial f/\partial x_{k}=-4k\,x_{k}^{-5}=-4k/x_{k}^{5}$.
  \item \textbf{The sum inside a function.} For $f=\bigl(\sum_{i}x_{i}\bigr)^{2}$
        use the chain rule; the inner derivative is
        $\partial\bigl(\sum_{i}x_{i}\bigr)/\partial x_{k}=1$, so
        $\partial f/\partial x_{k}=2\sum_{i=1}^{n}x_{i}$, the same for
        every $k$.
  \item \textbf{Products.} $f=\prod_{i}x_{i}$ gives
        $\partial f/\partial x_{k}=\prod_{i\neq k}x_{i}$: the product with
        the factor $x_{k}$ left out.
\end{itemize}
The summation index is a dummy and the differentiation variable is fixed:
call them by different letters ($i$ and $k$). Writing
$\partial/\partial x_{i}$ of a sum over $i$ is the standard source of
confusion.
\end{rulebox}

\begin{rulebox}{Exponentials to another base}
For $a>0$, $a^{x}=\mathrm{e}^{x\ln a}$, so
\[
  \frac{\mathrm{d}}{\mathrm{d}x}\,a^{x}=a^{x}\ln a,
  \qquad
  \pd{}{x}\,a^{g(x,y)}=a^{g(x,y)}\ln a\cdot g_{x}(x,y) .
\]
For example $\partial(2^{xy})/\partial x=y\,2^{xy}\ln2$. Keep the three
shapes apart: $x^{a}$ (variable base, power rule), $a^{x}$ (variable
exponent, factor $\ln a$), and $x^{y}$ (both, written as
$\mathrm{e}^{y\ln x}$).
\end{rulebox}

\begin{gallery}
\begin{tikzpicture}
  \begin{axis}[surfaxis, title={$f(K,L)=\sqrt{KL}$, sliced at $L=2$},
      xlabel={$K$}, ylabel={$L$}, zlabel={}, view={135}{26},
      zmin=0, zmax=3.2, z post scale=1.7]
    \addplot3[surf, domain=0.15:3, y domain=0.15:3, opacity=0.85]{sqrt(x*y)};
    \addplot3[samples=80, samples y=0, domain=0.15:3, line width=1.1pt,
      draw=black]({x},{2},{sqrt(2*x)});
  \end{axis}
\end{tikzpicture}\hfill
\begin{tikzpicture}
  \begin{axis}[contouraxis, xmin=-0.3, xmax=3.4, ymin=-0.3, ymax=3.4,
      xtick={1,2,3}, ytick={1,2,3},
      title={isoquants of $\sqrt{KL}$ with $\grad f$},
      xlabel={$K$}, ylabel={$L$}]
    \addplot[accent!40, line width=0.5pt, domain=0.22:3.3]{0.64/x};
    \addplot[accent!40, line width=0.5pt, domain=0.48:3.3]{1.44/x};
    \addplot[accent!40, line width=0.5pt, domain=0.85:3.3]{2.56/x};
    \addplot[accent!40, line width=0.5pt, domain=1.34:3.3]{4/x};
    \draw[-{Stealth[length=6pt,width=4pt]}, accent, line width=1pt]
      (axis cs:1,1) -- (axis cs:1.62,1.62);
    \draw[-{Stealth[length=6pt,width=4pt]}, accent, line width=1pt]
      (axis cs:2,0.5) -- (axis cs:2.16,1.14);
    \draw[-{Stealth[length=6pt,width=4pt]}, accent, line width=1pt]
      (axis cs:0.5,2) -- (axis cs:1.14,2.16);
    \addplot[only marks, mark=*, mark size=1pt, black, forget plot]
      coordinates {(1,1) (2,0.5) (0.5,2)};
  \end{axis}
\end{tikzpicture}
\gcap{Left: the plane $L=2$ cuts the surface; the slope of the cut curve is
$\partial f/\partial K$. Right: the gradient at three points, each one
perpendicular to the isoquant through it.}
\end{gallery}

\begin{rulebox}{Geometric meaning --- and why}
All three facts follow from one identity: the rate of change of $f$ in the
direction of a unit vector $\mathbf{u}$ is $\grad f\cdot\mathbf{u}$ (the
directional derivative below).
\begin{itemize}[leftmargin=1.6em,itemsep=0.4ex]
  \item \textbf{Steepest ascent.} Since
        $\grad f\cdot\mathbf{u}=\|\grad f\|\,\|\mathbf{u}\|\cos\theta
        =\|\grad f\|\cos\theta$, the rate is largest when $\cos\theta=1$,
        i.e. when $\mathbf{u}$ points \emph{along} $\grad f$. So
        $\grad f(\mathbf{x}_0)$ is the direction of steepest ascent, and
        $-\grad f$ of steepest descent.
  \item \textbf{Its length is that steepest rate.} Putting $\theta=0$ above
        leaves exactly $\|\grad f(\mathbf{x}_0)\|$.
  \item \textbf{Perpendicular to the level set.} Let $\mathbf{c}(t)$ be any
        curve running inside the level set through $\mathbf{x}_{0}$, so that
        $f(\mathbf{c}(t))$ is \emph{constant}. Differentiating with the
        chain rule gives
        \[
          0=\frac{\mathrm{d}}{\mathrm{d}t}f(\mathbf{c}(t))
           =\grad f(\mathbf{c}(t))\cdot\mathbf{c}'(t) .
        \]
        At $t_{0}$ this says $\grad f(\mathbf{x}_{0})$ is orthogonal to
        $\mathbf{c}'(t_{0})$, the tangent to the level set. Since the curve
        was arbitrary, the gradient is perpendicular to the whole level set.
\end{itemize}
The third fact is the one that matters most later: an optimum on a
constraint is a point where the level set of the objective and the
constraint set are \emph{tangent}, which is the same as their gradients
being parallel --- and that is the Lagrange condition.
\end{rulebox}

\begin{rulebox}{Economic interpretation}
\begin{itemize}[leftmargin=1.6em,itemsep=0.3ex]
  \item For a production function $Q=F(K,L)$: $F_K=\mathrm{MP}_K$
        (marginal product of capital), $F_L=\mathrm{MP}_L$ (marginal
        product of labour).
  \item For a utility function $U(x,y)$: $U_x, U_y$ are the marginal
        utilities; the marginal rate of substitution is
        $\mathrm{MRS}_{xy}=U_x/U_y$.
  \item For production: the marginal rate of technical substitution is
        $\mathrm{MRTS}_{LK}=F_L/F_K$.
\end{itemize}
\end{rulebox}

\begin{rulebox}{Directional derivative}
Let $\mathbf{u}\in\R^n$ be a \emph{unit} vector ($\|\mathbf{u}\|=1$). The
directional derivative of $f$ at $\mathbf{x}_0$ in the direction
$\mathbf{u}$ is
\[
  D_{\mathbf{u}}f(\mathbf{x}_0) \;=\; \grad f(\mathbf{x}_0)\cdot\mathbf{u} .
\]
\end{rulebox}

\begin{rulebox}{Caution}
If you are given a non-unit direction $\mathbf{v}$, you \emph{must}
normalise first: $\mathbf{u}=\mathbf{v}/\|\mathbf{v}\|$. Forgetting this is
the single most common error in directional-derivative problems.
\end{rulebox}

\begin{rulebox}{Differentiability in several variables}
Let $D\subseteq\R^{n}$ be open and $\mathbf{x}_{0}\in D$. The function $f$
is \textbf{differentiable} \de{(differenzierbar)} at $\mathbf{x}_{0}$ if
there is a vector $\mathbf{a}$ with
\[
  f(\mathbf{x}_{0}+\mathbf{h})=f(\mathbf{x}_{0})+\mathbf{a}\cdot\mathbf{h}
  +r(\mathbf{h}),
  \qquad
  \frac{r(\mathbf{h})}{\|\mathbf{h}\|}\to0\ \text{ as }\ \mathbf{h}\to\mathbf{0}.
\]
Near $\mathbf{x}_{0}$, $f$ is a linear function plus an error that vanishes
faster than the distance. For $n=2$ the graph of the linear part is a plane,
so differentiability says that the surface has a \textbf{tangent plane} at
the point: zoom in far enough and the surface looks flat.

If $f$ is differentiable at $\mathbf{x}_{0}$, then
\begin{itemize}[leftmargin=1.6em,itemsep=0.3ex]
  \item $f$ is continuous at $\mathbf{x}_{0}$;
  \item every partial derivative exists there, and
        $\mathbf{a}=\grad f(\mathbf{x}_{0})$ (take
        $\mathbf{h}=t\mathbf{e}_{k}$ and divide by $t$);
  \item every directional derivative exists and equals
        $\grad f(\mathbf{x}_{0})\cdot\mathbf{u}$.
\end{itemize}
\textbf{Sufficient criterion.} If all partials exist on an open
neighbourhood of $\mathbf{x}_{0}$ and are continuous there ($f$ is
$C^{1}$, \emph{continuously differentiable}), then $f$ is differentiable at
$\mathbf{x}_{0}$. Polynomials, exponentials, logarithms of positive
expressions, roots of positive expressions, quotients with non-zero
denominator, and all their sums, products and compositions are $C^{1}$ on
the open set where their formulas apply.

\textbf{Partials alone are not enough.} Take
$f(x,y)=\dfrac{xy}{x^{2}+y^{2}}$ with $f(0,0)=0$. On both axes $f=0$, so
$f_{x}(0,0)=f_{y}(0,0)=0$. But on the diagonal $y=x$,
$f=x^{2}/(2x^{2})=\tfrac12$ for every $x\neq0$. So $f$ is not even
continuous at the origin, and therefore not differentiable: the partials
see only two directions out of infinitely many.
\end{rulebox}

\begin{rulebox}{Where differentiability fails, and finding the set}
Compute the partials by the rules. The formulas are valid on the open set
where the rules applied; the points where a formula breaks down (a zero
denominator, a zero radicand, the kink of an absolute value or a minimum)
are the suspects, and each suspect is settled by the difference quotient.
\begin{itemize}[leftmargin=1.6em,itemsep=0.3ex]
  \item \textbf{Cone apex.} $f=\sqrt{(x-a)^{2}+(y-b)^{2}}$ has
        $f_{x}=(x-a)/f$, undefined at $(a,b)$. Directly,
        $f(a+h,b)-f(a,b)=\sqrt{h^{2}}=|h|$, and $|h|/h=\pm1$ has no limit:
        $f_{x}(a,b)$ does not exist, nor does $f_{y}(a,b)$. The same
        happens for the elliptic cone $\sqrt{4(x-a)^{2}+y^{2}}$ at $(a,0)$,
        with quotient $2|h|/h$. For the root of a sum of squares of
        \emph{linear} expressions, such as these cones, $f$ is
        differentiable exactly where the radicand is non-zero, and the
        graph has a sharp tip (or crease) where it vanishes. With higher
        powers under the root this can change: $\sqrt{x^{4}+y^{4}}$ is
        differentiable even at the origin, where its radicand is $0$.
  \item \textbf{$|x|+|y|$.} $f_{x}=1$ for $x>0$ and $-1$ for $x<0$; at
        $x=0$ the quotient is $|h|/h$ again. So $f_{x}$ fails on the whole
        $y$-axis and $f_{y}$ on the whole $x$-axis, and $f$ is
        differentiable exactly off the axes. The graph is an upside-down
        square pyramid, creased along both axes.
  \item \textbf{$\min\{ax,by\}$} with $a,b>0$. Where $ax<by$, $f=ax$;
        where $ax>by$, $f=by$. On the kink line $ax=by$ (in the quadrant
        $x,y>0$ of economics, a ray from the origin), raising $x$ leaves
        $f$ at $by$ (slope $0$) while lowering $x$ lowers $f$ at rate $a$:
        the one-sided slopes differ, so $f_{x}$ does not exist there, and
        nor does $f_{y}$. The graph is a roof whose ridge lies above that
        line.
\end{itemize}
A breaking formula is a warning, not a verdict: the partial may still
exist at that point, and only the difference quotient decides. And partials
that exist at a point still do not guarantee differentiability; continuity
of the partials near the point does.
\end{rulebox}

\begin{rulebox}{Proving differentiability from the definition}
When the $C^{1}$ criterion is not available, or a proof is asked for,
work with the remainder directly. At a point $\mathbf{x}^{*}$:
\begin{enumerate}[leftmargin=1.6em,itemsep=0.3ex]
  \item \textbf{Candidate.} Compute $\grad f(\mathbf{x}^{*})$; it is the
        only possible linear part. If a partial does not exist, stop: $f$
        is not differentiable.
  \item \textbf{Remainder.}
        $r(\mathbf{h})=f(\mathbf{x}^{*}+\mathbf{h})-f(\mathbf{x}^{*})
        -\grad f(\mathbf{x}^{*})\cdot\mathbf{h}$, expanded and simplified.
        For a polynomial, every term left over contains at least two
        factors $h_{i}$.
  \item \textbf{Estimate.} Show $|r(\mathbf{h})|/\|\mathbf{h}\|\to0$ with
        the standard inequalities
        \[
          |h_{i}|\le\|\mathbf{h}\|,\qquad
          \frac{h_{i}^{2}}{\|\mathbf{h}\|}\le\|\mathbf{h}\|,\qquad
          \frac{|h_{1}h_{2}|}{\|\mathbf{h}\|}\le\tfrac12\|\mathbf{h}\| ,
        \]
        and, for a product of three factors $h_{i}$, a bound
        $\|\mathbf{h}\|^{2}$.
  \item \textbf{Transcendental factors.} Replace them by their linear part
        plus a remainder: $\mathrm{e}^{t}=1+t+R(t)$ with $|R(t)|\le2t^{2}$
        for $|t|\le1$, and $\ln(1+t)=t+R(t)$ with $|R(t)|\le2t^{2}$ for
        $|t|\le\tfrac12$ (both from the Lagrange remainder of Chapter~4).
        The quadratic bound on $R$ is what makes $R/\|\mathbf{h}\|\to0$.
\end{enumerate}
\textbf{Roots at the origin.} For $f=(x^{2}+y^{2})^{p}$ with $p>0$,
$f(\mathbf{h})=\|\mathbf{h}\|^{2p}$. The difference quotient on an axis is
$|h|^{2p}/h$: it tends to $0$ if $2p>1$, and has no limit if $2p\le1$. So
$f$ is differentiable at the origin exactly when $p>\tfrac12$; the cone
$p=\tfrac12$ is the borderline case, and $p=\tfrac13$ is worse than a cone.
\end{rulebox}

\begin{rulebox}{Tangent plane and linearisation}
If $f$ is differentiable at $P=(x_{0},y_{0})$, the \textbf{tangent plane}
\de{(Tangentialebene)} to its graph at $(x_{0},y_{0},f(P))$ is
\[
  z=f(P)+f_{x}(P)\,(x-x_{0})+f_{y}(P)\,(y-y_{0}) .
\]
In $n$ variables, with $P=\mathbf{p}=(p_{1},\ldots,p_{n})$, the
\textbf{tangent hyperplane} is
\[
  z=f(\mathbf{p})+\grad f(\mathbf{p})\cdot(\mathbf{x}-\mathbf{p})
   =f(\mathbf{p})+\sum_{i=1}^{n}f_{x_i}(\mathbf{p})\,(x_{i}-p_{i}) .
\]
\begin{itemize}[leftmargin=1.6em,itemsep=0.3ex]
  \item The right-hand side is the \textbf{linearisation} $L(\mathbf{x})$
        of $f$ at $\mathbf{p}$, the multivariable version of the tangent
        line $f(a)+f'(a)(x-a)$ of Chapter~6. With
        $\mathrm{d}x_{i}=x_{i}-p_{i}$ it reads $L=f(\mathbf{p})+\mathrm{d}f$:
        the old value plus the total differential.
  \item The partials are evaluated \emph{at the point} $\mathbf{p}$; they
        are numbers, which is what makes the result a plane.
  \item For a point $Q$ near $P$, $f(Q)\approx L(Q)$. The difference
        $f(Q)-L(Q)$ is the remainder $r(\mathbf{h})$ of the definition, so
        comparing the two numbers measures how good the approximation is.
        For $f$ with continuous second partials it is at most of the order
        $\|Q-P\|^{2}$: halve the step and the error falls to about a
        quarter, or even less.
  \item Equivalently, $f_{x}(P)(x-x_{0})+f_{y}(P)(y-y_{0})-(z-f(P))=0$: the
        vector $(f_{x}(P),\,f_{y}(P),\,-1)$ is normal to the plane.
\end{itemize}
Where $f$ is not differentiable there is no tangent plane: a cone has none
at its apex.
\end{rulebox}

\section{Worked Examples}

\begin{worked}{For $f(x,y)=x^{3}y^{2}+\mathrm{e}^{2x}\ln y$, compute both partials and the gradient at $(1,1)$; then read the marginal products of $Q=10K^{0.3}L^{0.7}$}
\wstep{Step 1. Differentiating in $x$ means freezing $y$ --- literally.}
Treat $y$ as if it were the number $7$. Then $x^{3}y^{2}$ is
``$y^{2}$ times $x^{3}$'', a constant times a power, and $\ln y$ is a
constant multiplying $\mathrm{e}^{2x}$:
\[
  f_{x}=3x^{2}y^{2}+2\mathrm{e}^{2x}\ln y .
\]
\wtrap{Freezing means the frozen variable never produces a derivative
factor. $\frac{\partial}{\partial x}\bigl[\ln y\bigr]=0$, not $1/y$. The
symbol $\partial$ instead of $\mathrm{d}$ is the reminder that something is
being held fixed.}

\wstep{Step 2. Now freeze $x$ instead.}
$x^{3}$ is a constant multiplying $y^{2}$, and $\mathrm{e}^{2x}$ is a
constant multiplying $\ln y$:
\[
  f_{y}=2x^{3}y+\frac{\mathrm{e}^{2x}}{y} .
\]

\wstep{Step 3. Evaluate, and assemble the gradient.}
At $(1,1)$, with $\ln1=0$ and $\mathrm{e}^{2}\approx7.389$:
\[
  f_{x}(1,1)=3\cdot1\cdot1+2\mathrm{e}^{2}\cdot0=3,
  \qquad
  f_{y}(1,1)=2\cdot1\cdot1+\frac{\mathrm{e}^{2}}{1}=2+\mathrm{e}^{2}\approx9.389 .
\]
\[
  \grad f(1,1)=\begin{pmatrix}3\\2+\mathrm{e}^{2}\end{pmatrix}\approx\begin{pmatrix}3\\9.389\end{pmatrix}
\]
\wtrap{The gradient is a \textbf{vector}, and it is a vector in the
\emph{input} space --- here a direction in the $xy$-plane, not on the
surface. Reporting a single number, or a three-component object including
the height, are the two standard confusions.}

\wstep{Step 4. What each partial measures.}
$f_{x}$ is the slope of the curve obtained by slicing the surface with the
plane $y=1$ --- an ordinary one-variable derivative of the function
$x\mapsto f(x,1)$. Indeed $f(x,1)=x^{3}$, whose derivative at $x=1$ is $3$.
\checkmark\ Two partials, two slices, two ordinary derivatives.

\wstep{Step 5. What the gradient does that the partials alone do not.}
\begin{itemize}[leftmargin=1.6em,itemsep=0.2ex]
  \item It points in the direction of \textbf{steepest ascent}. Here that
        is roughly $(0.30,\,0.95)$ after normalising: mostly ``north''.
  \item Its length $\|\grad f\|=\sqrt{9+(2+\mathrm{e}^{2})^{2}}\approx9.86$
        is the steepest rate of climb available at that point.
  \item It is \textbf{perpendicular} to the level curve through the point.
\end{itemize}
The last is why the gradient governs constrained optimisation: at an
optimum the objective's gradient must be parallel to the constraint's, which
is the Lagrange condition.

\wstep{Step 6. Directional derivative: normalise, then take an inner product.}
In the direction of $\mathbf v=(3,4)$:
\[
  \mathbf u=\frac{\mathbf v}{\|\mathbf v\|}=\Bigl(\tfrac35,\tfrac45\Bigr),
  \qquad
  D_{\mathbf u}f=\grad f\cdot\mathbf u
  =\tfrac35\cdot3+\tfrac45(2+\mathrm{e}^{2})\approx9.31 .
\]
\wtrap{Normalise first. Using $(3,4)$ unnormalised gives $46.6$, which is
five times too large and has no meaning --- ``rate of change per unit
distance'' requires a unit vector. And $9.31<9.86$, as it must be: no
direction beats the gradient.}

\wstep{Step 7. The economic reading, on a production function.}
For $Q=10K^{0.3}L^{0.7}$:
\[
  MP_{K}=\pd{Q}{K}=3K^{-0.7}L^{0.7}=3\Bigl(\frac{L}{K}\Bigr)^{0.7},
  \qquad
  MP_{L}=\pd{Q}{L}=7\Bigl(\frac{K}{L}\Bigr)^{0.3} .
\]
At $K=100$, $L=50$: $MP_{K}=3(0.5)^{0.7}\approx1.85$ and
$MP_{L}=7(2)^{0.3}\approx8.62$. So one more unit of labour adds about $8.6$
units of output, one more of capital about $1.85$.

Two structural facts fall out. The \textbf{marginal rate of technical
substitution} is
\[
  \mathrm{MRTS}=\frac{MP_{L}}{MP_{K}}=\frac{0.7}{0.3}\cdot\frac{K}{L}\approx4.67 ,
\]
and each marginal product depends only on the \emph{ratio} $K/L$ --- the
signature of constant returns to scale, which is why Cobb--Douglas with
exponents summing to $1$ behaves the way it does.
\wtrap{$MP_{K}$ falls as $K$ rises (the exponent $-0.7$ is negative):
diminishing marginal product. But that is diminishing \emph{marginal},
not diminishing total --- output still increases. Confusing the two is the
most common misreading of a partial derivative in economics.}
\end{worked}

\begin{worked}{Find the tangent plane of $f(x,y)=x\ln(x+2y)$ at $P=(1,0)$, and test it at the nearby point $Q=(1.1,\,0.05)$}
\wstep{Step 1. Make sure there is a plane to find.}
The domain is the open half-plane $x+2y>0$, and $P$ lies in it since
$1+0>0$. On that half-plane $f$ is a polynomial times the logarithm of a
positive linear expression, so its partials exist and are continuous: $f$
is $C^{1}$, hence differentiable at $P$, and the tangent plane exists.

\wstep{Step 2. The height at $P$.}
$f(1,0)=1\cdot\ln1=0$. The plane touches the surface at $(1,0,0)$.

\wstep{Step 3. The two partials, then their values at $P$.}
In $x$ the function is a product of $x$ and $\ln(x+2y)$:
\[
  f_{x}=\ln(x+2y)+x\cdot\frac{1}{x+2y} .
\]
In $y$ the factor $x$ is a constant, and the chain rule supplies the inner
derivative $2$:
\[
  f_{y}=\frac{2x}{x+2y} .
\]
At $P$: $f_{x}(1,0)=\ln1+\tfrac11=1$ and $f_{y}(1,0)=\tfrac21=2$.
\wtrap{$x$ appears twice in $f$, once as a factor and once inside the
logarithm, so $f_{x}$ needs the product rule. Writing $f_{x}=x/(x+2y)$
alone loses the term $\ln(x+2y)$; here it happens to vanish at $P$, which
hides the error until the next point.}

\wstep{Step 4. Assemble the plane.}
\[
  z=f(P)+f_{x}(P)(x-1)+f_{y}(P)(y-0)=0+1\cdot(x-1)+2y=x+2y-1 .
\]
\wtrap{The partials go in as \emph{numbers} evaluated at $P$. Leaving them
as functions of $x$ and $y$ gives a curved surface, not a plane.}

\wstep{Step 5. Compare at $Q$.}
\[
  f(Q)=1.1\cdot\ln(1.1+0.1)=1.1\ln1.2\approx0.200554,
  \qquad
  L(Q)=1.1+0.1-1=0.200000 .
\]
The plane misses by about $0.000554$, while $f$ itself changed by about
$0.2$: the error is under $0.3\%$ of the change. The step has length
$\|Q-P\|=\sqrt{0.1^{2}+0.05^{2}}\approx0.112$. The error is far below that
length, and even below its square $0.0125$, as the rule box predicts: the
remainder vanishes faster than the distance.

\wstep{Step 6. The same number as a total differential.}
With $\mathrm{d}x=0.1$ and $\mathrm{d}y=0.05$,
\[
  \mathrm{d}f=f_{x}(P)\,\mathrm{d}x+f_{y}(P)\,\mathrm{d}y=0.1+0.1=0.2,
\]
and $L(Q)=f(P)+\mathrm{d}f=0.2$. \checkmark\ The tangent plane and the
total differential are one idea written two ways.
\end{worked}

\begin{worked}{Prove from the definition that $f(x,y)=x^{2}y$ is differentiable at every point; then decide at the origin for $(x^{2}+y^{2})^{1/3}$ and $(x^{2}+y^{2})^{2/3}$}
\wstep{Step 1. The candidate for the linear part.}
Fix a point $(a,b)$. If $f$ is differentiable there, the linear part must
be $\grad f(a,b)\cdot\mathbf{h}$, so compute the partials:
$f_{x}=2xy$, $f_{y}=x^{2}$, hence $\grad f(a,b)=(2ab,\,a^{2})$.

\wstep{Step 2. Form the remainder and let everything cancel that can.}
With $\mathbf{h}=(h_{1},h_{2})$,
\[
  f(a+h_{1},b+h_{2})=(a^{2}+2ah_{1}+h_{1}^{2})(b+h_{2})
  =a^{2}b+a^{2}h_{2}+2abh_{1}+2ah_{1}h_{2}+bh_{1}^{2}+h_{1}^{2}h_{2} .
\]
Subtract $f(a,b)=a^{2}b$ and the linear part $2abh_{1}+a^{2}h_{2}$:
\[
  r(\mathbf{h})=2a\,h_{1}h_{2}+b\,h_{1}^{2}+h_{1}^{2}h_{2} .
\]
Every surviving term contains at least \emph{two} factors $h_{i}$. That is
the sign that the candidate was right.

\wstep{Step 3. Bound the quotient by something that tends to $0$.}
Use $|h_{i}|\le\|\mathbf{h}\|$ and $|h_{1}h_{2}|\le\tfrac12\|\mathbf{h}\|^{2}$
(from $(|h_{1}|-|h_{2}|)^{2}\ge0$):
\[
  \frac{|r(\mathbf{h})|}{\|\mathbf{h}\|}
  \le\frac{2|a|\cdot\tfrac12\|\mathbf{h}\|^{2}+|b|\,\|\mathbf{h}\|^{2}
          +\|\mathbf{h}\|^{3}}{\|\mathbf{h}\|}
  =(|a|+|b|)\,\|\mathbf{h}\|+\|\mathbf{h}\|^{2}\;\longrightarrow\;0 .
\]
So $f$ is differentiable at $(a,b)$, for every $(a,b)$, with derivative
$(2ab,\,a^{2})$.
\wtrap{Checking $r/\|\mathbf{h}\|\to0$ along the axes, or along every
line, is not a proof: the bound must hold for \emph{all} small $\mathbf{h}$
at once. The bound may depend on the fixed point $(a,b)$, but only through
constants.}

\wstep{Step 4. A root at the origin: $g=(x^{2}+y^{2})^{1/3}$.}
First the partials, from the difference quotient:
\[
  \frac{g(h,0)-g(0,0)}{h}=\frac{(h^{2})^{1/3}}{h}
  =\frac{|h|^{2/3}}{h}=\pm|h|^{-1/3},
\]
which is unbounded as $h\to0$. So $g_{x}(0,0)$ does not exist (nor, by
symmetry, $g_{y}(0,0)$), and $g$ is \textbf{not differentiable} at the
origin, although it is continuous there. Its graph has a tip even sharper
than a cone's.

\wstep{Step 5. A higher power: $k=(x^{2}+y^{2})^{2/3}$.}
Now $\dfrac{k(h,0)-k(0,0)}{h}=\dfrac{|h|^{4/3}}{h}=\pm|h|^{1/3}\to0$, so both
partials are $0$ and the candidate linear part is $0$. The remainder is
$r(\mathbf{h})=k(\mathbf{h})=\|\mathbf{h}\|^{4/3}$, and
\[
  \frac{|r(\mathbf{h})|}{\|\mathbf{h}\|}=\|\mathbf{h}\|^{1/3}\to0 .
\]
So $k$ \textbf{is} differentiable at the origin, with the tangent plane
$z=0$. In general $(x^{2}+y^{2})^{p}=\|\mathbf{h}\|^{2p}$ at the origin, and
the verdict depends only on whether $2p>1$.
\end{worked}

\section{Practice Problems}

\begin{problems}
  \item Compute $f_x$ and $f_y$ for $f(x,y)=5x^4-2x^2y^3+y^5$.
  \item For $F(K,L)=12\,K^{1/4}L^{3/4}$, find $F_K$ and $F_L$ at
        $(K,L)=(16,81)$.
  \item Compute $f_x$ and $f_y$ for $f(x,y)=x^2e^{3y}$.
  \item Compute $\nabla f$ for $f(x,y)=4x^3y-xy^4+7$.
  \item Compute the directional derivative of $f(x,y)=x^2-4xy+y^3$ at
        $(2,1)$ in the direction $\mathbf{v}=(5,-12)$.
  \item For $F(K,L)=6K^{2/3}L^{1/3}$, compute $\mathrm{MP}_K$ and
        $\mathrm{MP}_L$ at $(K,L)=(27,8)$.
  \item Compute $f_x$ and $f_y$ for $f(x,y)=(3x-2y)^4$.
  \item Compute $f_x$ and $f_y$ for $f(x,y)=e^{xy^2}$.
  \item For $F(K,L)=25K^{0.6}L^{0.5}$, find $F_K$ and $F_L$.
  \item Compute $\nabla f$ for $f(x,y,z)=x^2+3yz-z^2e^x$.
  \item Compute the directional derivative of $f(x,y)=e^x\ln(1+y)$ at
        $(0,0)$ in the direction $\mathbf{v}=(1,2)$.
  \item For $U(x,y)=2\ln(x)+5\ln(y)$, derive $\mathrm{MRS}_{xy}=U_x/U_y$ as
        a function of $(x,y)$ and evaluate at $(3,5)$.
  \item Simplify and compute $f_x, f_y$ for
        $f(x,y)=\dfrac{x^2-y^2}{x+y}$.
  \item Compute $f_x$ and $f_y$ for $f(x,y)=\ln(x^2y+1)$.
  \item For $U(x_1,x_2,x_3)=x_1^{1/2}x_2^{1/3}x_3^{1/6}$, compute the
        three partials.
  \item Compute the directional derivative of $f(x,y,z)=xy-2yz+z^2$ at
        $(1,2,3)$ in the direction $\mathbf{v}=(1,2,2)$.
  \item For $F(K,L)=K^{0.25}L^{0.75}$, derive $\mathrm{MRTS}_{LK}=F_L/F_K$
        and evaluate at $(K,L)=(16,4)$.
  \item Expand and then compute $f_x, f_y$ for
        $f(x,y)=(x^2+y)(x-3y^2)$.
  \item Compute $f_x$ and $f_y$ for $f(x,y)=xe^{-y/x}$ with $x\neq0$.
  \item For $F(K,L)=(3K^{1/2}+2L^{1/2})^2$, find $F_K$ and $F_L$.
  \item Compute $\nabla f$ for $f(x,y)=\arctan(y/x)$ with $x>0$.
  \item For $f(x,y)=x^{1/2}y^{1/2}$ at $(K,L)=(4,9)$, find the unit
        direction of steepest ascent and the maximum rate of increase.
  \item For the CES function $F(K,L)=(K^{1/2}+L^{1/2})^2$, verify Euler's
        theorem $KF_K+LF_L=rF$ for the appropriate degree of homogeneity
        $r$.
  \item Compute $f_x$ and $f_y$ for $f(x,y)=\dfrac{2x+3y}{x^2+y^2}$.
  \item Compute $f_x$ and $f_y$ for $f(x,y)=\dfrac{e^x}{1+e^y}$.
  \item For $F(K,L)=A(\alpha K^\rho+(1-\alpha)L^\rho)^{\nu/\rho}$ (CES),
        find $F_L$.
  \item Compute the directional derivative of $f(x,y)=\ln(x^2+y^2)$ at
        $(3,4)$ in the direction $\mathbf{v}=(-1,1)$.
  \item Compute $f_x, f_y, f_z$ for $f(x,y,z)=x^2yz^3-4xy^2z+5z$.
  \item Compute $f_x$ and $f_y$ for $f(x,y)=\sqrt{4x^2+9y^2}$.
  \item For $F(K,L)=(K^{1/2}+L^{1/2})^4$, compute $F_K$ at
        $(K,L)=(4,9)$.
  \item Compute $\nabla f$ for $f(x,y,z)=\ln(x^2+y^2+z^2)$.
  \item Compute $f_x$ and $f_y$ for
        $f(x,y)=\ln\left(\sqrt{x^2+y^2}\right)$.
  \item Cost minimisation. A firm uses inputs $(x,y)$ with cost
        $C(x,y)=4x+9y$ and output $Q(x,y)=x^{1/2}y^{1/2}$. Compute
        $\nabla C$ and $\nabla Q$, and verify that at the cost-minimising
        input mix $\nabla C=\lambda\nabla Q$ for some $\lambda$.
  \item Compute $f_x$ and $f_y$ for $f(x,y)=(1+xy)^5$.
  \item Compute the directional derivative of $f(x,y)=x^3-2xy+y^2$ at
        $(1,2)$ in the direction $\mathbf{v}=(3,4)$.
  \item For $F(K,L)=10K^{0.3}L^{0.7}$ (Cobb--Douglas), compute $F_K$ and
        $F_L$ at $(K,L)=(100,100)$ and verify $KF_K+LF_L=F$.
  \item Compute $f_x$ and $f_y$ for $f(x,y)=\sin(xy)+\cos(x-y)$.
  \item Compute $\nabla f$ for $f(x,y)=x^2y-xy^2+3x-2y$ at the point
        $(1,-1)$.
  \item For $U(x,y)=x^\alpha y^{1-\alpha}$ (Cobb--Douglas,
        $0<\alpha<1$), derive $\mathrm{MRS}_{xy}=U_x/U_y$ and show it
        equals $\dfrac{\alpha}{1-\alpha}\cdot\dfrac{y}{x}$.
  \item Compute the directional derivative of $f(x,y)=e^{xy}$ at $(1,0)$
        in the direction $\mathbf{v}=(-3,4)$.
  \item Compute $f_x$ and $f_y$ for $f(x,y)=x^y$ (with $x>0$).
        \emph{Hint:} write $x^y=e^{y\ln x}$.
  \item For $F(K,L)=\min\{2K,3L\}$ (Leontief, fixed-proportions, on the
        smooth region away from the kink), compute $F_K$ and $F_L$ in
        each of the two regions $\{2K<3L\}$ and $\{2K>3L\}$.
  \item Compute $\nabla f$ for $f(x,y,z)=xyz+x^2-z^2$ at the point
        $(2,-1,3)$.
  \item Compute the unit vector of steepest \emph{descent} for
        $f(x,y)=x^2+4y^2$ at $(3,1)$, and the corresponding rate of
        decrease.
  \item For $f(x,y)=xe^y-ye^x$, compute $f_x, f_y$.
  \item Compute $f_x$ and $f_y$ for $f(x,y)=\dfrac{x}{x^2+y^2}$.
  \item For $U(x,y)=x+2\sqrt{y}$ (quasi-linear), derive
        $\mathrm{MRS}_{xy}$.
  \item Compute the directional derivative of $f(x,y,z)=x^2+y^2-z^2$ at
        $(2,-1,2)$ in the direction $\mathbf{v}=(2,2,1)$.
  \item Compute $f_x$ and $f_y$ for $f(x,y)=\tan(x+y)+\tan(x-y)$.
  \item For $F(K,L)=(\alpha K+\beta L)^r$ (linear-aggregate,
        $\alpha,\beta,r>0$), compute $F_K, F_L$, and verify
        $KF_K+LF_L=rF$.
  \item Compute $\nabla f$ for $f(x,y)=e^{-x}\sin(y)+e^x\cos(y)$.
  \item Compute the directional derivative of $f(x,y)=\sqrt{x^2+y^2}$ at
        $(3,4)$ in the direction $\mathbf{v}=(1,0)$.
  \item For $U(x,y)=\min\{2x,3y\}$, compute the marginal utilities on
        each smooth region $\{2x<3y\}$ and $\{2x>3y\}$.
  \item Compute $f_x, f_y$ for $f(x,y)=\dfrac{\ln(x+y)}{x-y}$ (assume
        $x>y>0$).
  \item For $F(K,L)=AK^\alpha L^\beta$ general Cobb--Douglas, derive
        $\mathrm{MRTS}_{LK}$ and show it depends only on the ratio $K/L$
        (not on scale).
  \item Compute the unit direction of steepest ascent for
        $f(x,y)=x^2-y^2$ at $(2,-1)$, and the maximum rate.
  \item Compute $f_x$ and $f_y$ for $f(x,y)=\cosh(x+y)+\sinh(x-y)$.
  \item For $f(x,y,z)=e^{x+y+z}$, compute $\nabla f$ at $(0,0,0)$ and the
        directional derivative in $\mathbf{v}=(1,1,1)$.
  \item Compute $f_x$ and $f_y$ for $f(x,y)=(\ln x)(\ln y)$ with
        $x,y>1$.
  \item For $F(K,L)=K+2K^{1/2}L^{1/2}+L$ (which equals
        $(K^{1/2}+L^{1/2})^2$), compute $F_K, F_L$ directly and confirm
        agreement with the CES form.
  \item Compute the directional derivative of $f(x,y)=x^2y+y^3$ at
        $(-1,2)$ in the direction $\mathbf{v}=(4,-3)$.
  \item For $U(x,y)=\alpha\ln x+\beta\ln y$ (log-Cobb--Douglas), derive
        $\mathrm{MRS}_{xy}$ and evaluate at $(x,y)=(\alpha,\beta)$.
  \item Compute $\nabla f$ for $f(x,y)=(x^2+y^2)^{-1/2}$ (Newtonian
        potential in 2D).
  \item Compute $f_x$ and $f_y$ for $f(x,y)=\arctan(xy)$.
  \item For $F(K,L)=(K^\rho+L^\rho)^{1/\rho}$ with $\rho<1$, $\rho\neq0$
        (CES with constant returns), compute $F_K$ and show
        $F_K=(K/F)^{\rho-1}$.
  \item Compute the directional derivative of $f(x,y,z)=xyz$ at
        $(2,3,-1)$ in the direction $\mathbf{v}=(2,-1,2)$.
\end{problems}

\section{Solutions}

\begin{solutions}
  \item $f_x=20x^3-4xy^3$, $f_y=-6x^2y^2+5y^4$.

  \item $F_K=12\cdot\frac14K^{-3/4}L^{3/4}=3(L/K)^{3/4}$;
        $F_L=12\cdot\frac34K^{1/4}L^{-1/4}=9(K/L)^{1/4}$. At $(16,81)$:
        $(L/K)^{3/4}=(81/16)^{3/4}=27/8$ so $F_K=81/8$;
        $(K/L)^{1/4}=(16/81)^{1/4}=2/3$ so $F_L=6$.

  \item $f_x=2xe^{3y}$, $f_y=3x^2e^{3y}$.

  \item $\nabla f=(12x^2y-y^4,\ 4x^3-4xy^3)^T$.

  \item $f_x=2x-4y$, $f_y=-4x+3y^2$. At $(2,1)$: $\nabla f=(0,-5)^T$.
        Normalise: $\|\mathbf{v}\|=\sqrt{25+144}=13$,
        $\mathbf{u}=(5/13,-12/13)^T$. Then
        \[
          D_{\mathbf{u}}f=0\cdot\frac{5}{13}+(-5)\cdot
            \left(-\frac{12}{13}\right)=\frac{60}{13} .
        \]

  \item $F_K=6\cdot\frac23K^{-1/3}L^{1/3}=4(L/K)^{1/3}$,
        $F_L=2(K/L)^{2/3}$. At $(27,8)$: $(L/K)^{1/3}=2/3$ so
        $\mathrm{MP}_K=8/3$; $(K/L)^{2/3}=(27/8)^{2/3}=9/4$ so
        $\mathrm{MP}_L=9/2$.

  \item Chain rule with $u=3x-2y$: $f_x=4(3x-2y)^3\cdot3=12(3x-2y)^3$,
        $f_y=-8(3x-2y)^3$.

  \item $f_x=y^2e^{xy^2}$, $f_y=2xye^{xy^2}$.

  \item $F_K=15K^{-0.4}L^{0.5}$, $F_L=12.5K^{0.6}L^{-0.5}$.

  \item $f_x=2x-z^2e^x$, $f_y=3z$, $f_z=3y-2ze^x$. So
        $\nabla f=(2x-z^2e^x,\ 3z,\ 3y-2ze^x)^T$.

  \item $f_x=e^x\ln(1+y)$, $f_y=e^x/(1+y)$. At $(0,0)$:
        $\nabla f=(0,1)^T$. Normalise: $\|\mathbf{v}\|=\sqrt5$,
        $\mathbf{u}=(1/\sqrt5,2/\sqrt5)^T$. So
        \[
          D_{\mathbf{u}}f=0\cdot\frac{1}{\sqrt5}+1\cdot\frac{2}{\sqrt5}
            =\frac{2\sqrt5}{5} .
        \]

  \item $U_x=2/x$, $U_y=5/y$. $\mathrm{MRS}_{xy}=U_x/U_y
        =(2/x)/(5/y)=2y/(5x)$. At $(3,5)$: $\mathrm{MRS}=10/15=2/3$.

  \item For $x+y\neq0$: $f(x,y)=x-y$, so $f_x=1$, $f_y=-1$.

  \item With $u=x^2y+1$: $f_x=2xy/(x^2y+1)$, $f_y=x^2/(x^2y+1)$.

  \item $U_{x_1}=\frac12x_1^{-1/2}x_2^{1/3}x_3^{1/6}$,
        $U_{x_2}=\frac13x_1^{1/2}x_2^{-2/3}x_3^{1/6}$,
        $U_{x_3}=\frac16x_1^{1/2}x_2^{1/3}x_3^{-5/6}$. Equivalently
        $U_{x_i}=(\alpha_i/x_i)U$ with $\alpha=(1/2,1/3,1/6)$.

  \item $f_x=y$, $f_y=x-2z$, $f_z=-2y+2z$. At $(1,2,3)$:
        $\nabla f=(2,-5,2)^T$. Normalise: $\|\mathbf{v}\|=\sqrt{1+4+4}=3$,
        $\mathbf{u}=(1/3,2/3,2/3)^T$. Then
        $D_{\mathbf{u}}f=(2-10+4)/3=-4/3$.

  \item $F_K=0.25K^{-0.75}L^{0.75}$, $F_L=0.75K^{0.25}L^{-0.25}$. So
        $\mathrm{MRTS}_{LK}=F_L/F_K=3K/L$. At $(16,4)$:
        $\mathrm{MRTS}=48/4=12$.

  \item Expand: $f=x^3-3x^2y^2+xy-3y^3$. Hence $f_x=3x^2-6xy^2+y$,
        $f_y=-6x^2y+x-9y^2$.

  \item With $u=-y/x$: $\partial u/\partial x=y/x^2$,
        $\partial u/\partial y=-1/x$. Product rule:
        \[
          f_x=e^{-y/x}+x\cdot e^{-y/x}\cdot\frac{y}{x^2}
            =e^{-y/x}\left(1+\frac{y}{x}\right), \qquad
          f_y=x\cdot e^{-y/x}\cdot\left(-\frac1x\right)=-e^{-y/x} .
        \]

  \item Let $g=3K^{1/2}+2L^{1/2}$, so $F=g^2$. Then
        \[
          F_K=2g\cdot\frac32K^{-1/2}=\frac{3g}{K^{1/2}}
            =\frac{3(3K^{1/2}+2L^{1/2})}{K^{1/2}}, \qquad
          F_L=2g\cdot L^{-1/2}=\frac{2(3K^{1/2}+2L^{1/2})}{L^{1/2}} .
        \]

  \item With $u=y/x$, $f=\arctan u$, $f'=1/(1+u^2)$:
        \[
          f_x=\frac{1}{1+y^2/x^2}\cdot\left(-\frac{y}{x^2}\right)
            =\frac{-y}{x^2+y^2}, \qquad
          f_y=\frac{x}{x^2+y^2} .
        \]
        So $\nabla f=(-y/(x^2+y^2),\ x/(x^2+y^2))^T$.

  \item $f_K=\frac12K^{-1/2}L^{1/2}$, $f_L=\frac12K^{1/2}L^{-1/2}$. At
        $(4,9)$: $\nabla f=(\frac12\cdot\frac32,\ \frac12\cdot\frac23)^T
        =(3/4,1/3)^T$. Then $\|\nabla f\|^2=9/16+1/9=97/144$, so the
        maximum rate is $\sqrt{97}/12\approx0.821$ in direction
        $\mathbf{u}^*=(9,4)^T/\sqrt{97}$.

  \item $F=(K^{1/2}+L^{1/2})^2$. Replacing $(K,L)\to(\lambda K,\lambda L)$
        gives $\lambda F$, so $r=1$. With $g=K^{1/2}+L^{1/2}$:
        $F_K=2g\cdot\frac12K^{-1/2}=g/K^{1/2}$, $F_L=g/L^{1/2}$. Hence
        \[
          KF_K+LF_L=K^{1/2}g+L^{1/2}g=g(K^{1/2}+L^{1/2})=g^2=F. \checkmark
        \]

  \item Quotient rule, $D=x^2+y^2$:
        \[
          f_x=\frac{2D-(2x+3y)(2x)}{D^2}=\frac{-2x^2-6xy+2y^2}{(x^2+y^2)^2},
          \qquad f_y=\frac{3x^2-4xy-3y^2}{(x^2+y^2)^2} .
        \]

  \item Treat denominator as constant for $f_x$: $f_x=e^x/(1+e^y)$. For
        $f_y$:
        \[
          f_y=e^x\cdot\frac{-e^y}{(1+e^y)^2}=-\frac{e^{x+y}}{(1+e^y)^2} .
        \]

  \item Let $S=\alpha K^\rho+(1-\alpha)L^\rho$, $F=AS^{\nu/\rho}$.
        Treating $K$ constant, $\partial S/\partial L=\rho(1-\alpha)L^{\rho-1}$:
        \[
          F_L=A\cdot\frac{\nu}{\rho}S^{\nu/\rho-1}\cdot\rho(1-\alpha)L^{\rho-1}
            =\nu(1-\alpha)AL^{\rho-1}S^{\nu/\rho-1} .
        \]

  \item $f=\ln(x^2+y^2)$, $f_x=2x/(x^2+y^2)$, $f_y=2y/(x^2+y^2)$. At
        $(3,4)$: $x^2+y^2=25$, $\nabla f=(6/25,8/25)^T$. Normalise:
        $\|\mathbf{v}\|=\sqrt2$, $\mathbf{u}=(-1/\sqrt2,1/\sqrt2)^T$. Then
        \[
          D_{\mathbf{u}}f=\frac{6}{25}\cdot\left(-\frac{1}{\sqrt2}\right)
            +\frac{8}{25}\cdot\frac{1}{\sqrt2}=\frac{2}{25\sqrt2}
            =\frac{\sqrt2}{25} .
        \]

  \item $f_x=2xyz^3-4y^2z$, $f_y=x^2z^3-8xyz$,
        $f_z=3x^2yz^2-4xy^2+5$.

  \item Chain rule, $u=4x^2+9y^2$:
        \[
          f_x=\frac{8x}{2\sqrt u}=\frac{4x}{\sqrt{4x^2+9y^2}}, \qquad
          f_y=\frac{9y}{\sqrt{4x^2+9y^2}} .
        \]

  \item With $g=K^{1/2}+L^{1/2}$, $F=g^4$,
        $F_K=4g^3\cdot\frac12K^{-1/2}=2g^3/K^{1/2}$. At $(4,9)$: $g=5$,
        $K^{1/2}=2$, so $F_K=2\cdot125/2=125$.

  \item Let $r^2=x^2+y^2+z^2$. Then $f=\ln(r^2)$ and $f_{x_i}=2x_i/r^2$.
        So $\nabla f=(2/(x^2+y^2+z^2))(x,y,z)^T$.

  \item Use $\ln\sqrt{u}=\frac12\ln u$: $f=\frac12\ln(x^2+y^2)$. So
        $f_x=x/(x^2+y^2)$, $f_y=y/(x^2+y^2)$.

  \item $\nabla C=(4,9)^T$. With $Q=x^{1/2}y^{1/2}$: $Q_x=\frac12(y/x)^{1/2}$
        and $Q_y=\frac12(x/y)^{1/2}$, so
        \[
          \nabla Q=\tfrac12\left((y/x)^{1/2},\,(x/y)^{1/2}\right)^{T} .
        \]
        The condition $\nabla C=\lambda\nabla Q$ requires
        $4/9=Q_x/Q_y=y/x$, i.e.\ $y=4x/9$. Along this locus
        $\nabla Q\propto(2/3,3/2)^T\propto(4,9)^T=\nabla C$. \checkmark

  \item Chain rule: $f_x=5(1+xy)^4\cdot y=5y(1+xy)^4$,
        $f_y=5x(1+xy)^4$.

  \item $f_x=3x^2-2y$, $f_y=-2x+2y$. At $(1,2)$:
        $\nabla f=(3-4,-2+4)^T=(-1,2)^T$. Normalise: $\|\mathbf{v}\|=5$,
        $\mathbf{u}=(3/5,4/5)^T$. Then
        \[
          D_{\mathbf{u}}f=-1\cdot\frac35+2\cdot\frac45
            =\frac{-3+8}{5}=1 .
        \]

  \item $F_K=10\cdot0.3K^{-0.7}L^{0.7}=3(L/K)^{0.7}$;
        $F_L=7(K/L)^{0.3}$. At $K=L=100$: $F_K=3$, $F_L=7$. And
        $F=10\cdot100^{0.3}\cdot100^{0.7}=10\cdot100=1000$. Then
        $KF_K+LF_L=100\cdot3+100\cdot7=1000=F$. \checkmark (homogeneity
        degree $r=0.3+0.7=1$).

  \item $f_x=y\cos(xy)-\sin(x-y)$, $f_y=x\cos(xy)+\sin(x-y)$.

  \item $f_x=2xy-y^2+3$, $f_y=x^2-2xy-2$. At $(1,-1)$: $f_x=-2-1+3=0$,
        $f_y=1+2-2=1$. So $\nabla f(1,-1)=(0,1)^T$.

  \item $U_x=\alpha x^{\alpha-1}y^{1-\alpha}$,
        $U_y=(1-\alpha)x^\alpha y^{-\alpha}$. So
        \[
          \mathrm{MRS}_{xy}=\frac{U_x}{U_y}
            =\frac{\alpha x^{\alpha-1}y^{1-\alpha}}{(1-\alpha)x^\alpha y^{-\alpha}}
            =\frac{\alpha}{1-\alpha}\cdot\frac{y}{x} . \checkmark
        \]

  \item $f_x=ye^{xy}$, $f_y=xe^{xy}$. At $(1,0)$: $\nabla f=(0,1)^T$.
        Normalise: $\|\mathbf{v}\|=5$, $\mathbf{u}=(-3/5,4/5)^T$. Then
        $D_{\mathbf{u}}f=0\cdot(-3/5)+1\cdot(4/5)=4/5$.

  \item Write $f=e^{y\ln x}$. Then $f_x=e^{y\ln x}\cdot y/x=yx^{y-1}$ and
        $f_y=e^{y\ln x}\cdot\ln x=x^y\ln x$.

  \item On $\{2K<3L\}$ the binding factor is capital: $F=2K$, so
        $F_K=2$, $F_L=0$. On $\{2K>3L\}$ labour binds: $F=3L$, so
        $F_K=0$, $F_L=3$. (Non-differentiable on $\{2K=3L\}$.)

  \item $f_x=yz+2x$, $f_y=xz$, $f_z=xy-2z$. At $(2,-1,3)$:
        $f_x=-3+4=1$, $f_y=6$, $f_z=-2-6=-8$. So
        $\nabla f=(1,6,-8)^T$.

  \item $\nabla f=(2x,8y)^T$. At $(3,1)$: $\nabla f=(6,8)^T$. Steepest
        descent is $-\nabla f/\|\nabla f\|$. $\|\nabla f\|=10$, so unit
        descent direction is $(-3/5,-4/5)^T$, with rate of decrease
        $-\|\nabla f\|=-10$ (i.e.\ $f$ decreases at rate 10).

  \item $f_x=e^y-ye^x$, $f_y=xe^y-e^x$.

  \item Quotient rule with $D=x^2+y^2$:
        \[
          f_x=\frac{D-x\cdot2x}{D^2}=\frac{y^2-x^2}{(x^2+y^2)^2}, \qquad
          f_y=\frac{-x\cdot2y}{D^2}=\frac{-2xy}{(x^2+y^2)^2} .
        \]

  \item $U_x=1$, $U_y=1/\sqrt{y}$. So $\mathrm{MRS}_{xy}=U_x/U_y
        =\sqrt{y}$ (depends only on $y$, quasi-linear utility).

  \item $f_x=2x$, $f_y=2y$, $f_z=-2z$. At $(2,-1,2)$:
        $\nabla f=(4,-2,-4)^T$. Normalise: $\|\mathbf{v}\|=\sqrt{4+4+1}=3$,
        $\mathbf{u}=(2/3,2/3,1/3)^T$. Then
        \[
          D_{\mathbf{u}}f=4\cdot\frac23+(-2)\cdot\frac23+(-4)\cdot\frac13
            =\frac{8-4-4}{3}=0 .
        \]
        (We are moving tangent to the level surface.)

  \item $f_x=\sec^2(x+y)+\sec^2(x-y)$, $f_y=\sec^2(x+y)-\sec^2(x-y)$.

  \item Let $S=\alpha K+\beta L$, $F=S^r$. Then $F_K=rS^{r-1}\alpha$,
        $F_L=rS^{r-1}\beta$. Hence
        \[
          KF_K+LF_L=rS^{r-1}(\alpha K+\beta L)=rS^{r-1}\cdot S=rS^r=rF.
            \checkmark
        \]

  \item $\nabla f=(-e^{-x}\sin y+e^x\cos y,\ e^{-x}\cos y-e^x\sin y)^T$.

  \item $f_x=x/\sqrt{x^2+y^2}$, $f_y=y/\sqrt{x^2+y^2}$. At $(3,4)$:
        $\sqrt{x^2+y^2}=5$, $\nabla f=(3/5,4/5)^T$. The direction
        $\mathbf{v}=(1,0)$ is already unit. So
        $D_{\mathbf{u}}f=3/5\cdot1+4/5\cdot0=3/5$.

  \item On $\{2x<3y\}$: $U=2x$, so $U_x=2$, $U_y=0$. On $\{2x>3y\}$:
        $U=3y$, so $U_x=0$, $U_y=3$. (Non-differentiable on the ray
        $2x=3y$.)

  \item Let $L=\ln(x+y)$, $D=x-y$. By the quotient rule:
        \[
          f_x=\frac{(1/(x+y))(x-y)-\ln(x+y)\cdot1}{(x-y)^2}
            =\frac{1/(x+y)\cdot(x-y)-\ln(x+y)}{(x-y)^2} .
        \]
        Simplifying
        $f_x=\dfrac{(x-y)/(x+y)-\ln(x+y)}{(x-y)^2}$. Similarly,
        \[
          f_y=\frac{(1/(x+y))(x-y)-\ln(x+y)\cdot(-1)}{(x-y)^2}
            =\frac{(x-y)/(x+y)+\ln(x+y)}{(x-y)^2} .
        \]

  \item $F_K=A\alpha K^{\alpha-1}L^\beta$, $F_L=A\beta K^\alpha L^{\beta-1}$.
        So
        \[
          \mathrm{MRTS}_{LK}=\frac{F_L}{F_K}=\frac{\beta}{\alpha}\cdot
            \frac{K}{L} .
        \]
        This depends only on the ratio $K/L$ and the exponent ratio
        $\beta/\alpha$: it is scale-invariant.

  \item $\nabla f=(2x,-2y)^T$. At $(2,-1)$: $\nabla f=(4,2)^T$. Direction
        of steepest ascent: $\mathbf{u}^*=(4,2)/\sqrt{20}=(2,1)/\sqrt5$.
        Maximum rate $=\|\nabla f\|=\sqrt{20}=2\sqrt5$.

  \item $f_x=\sinh(x+y)+\cosh(x-y)$, $f_y=\sinh(x+y)-\cosh(x-y)$.

  \item $\nabla f=(e^{x+y+z},e^{x+y+z},e^{x+y+z})^T$. At $(0,0,0)$:
        $\nabla f=(1,1,1)^T$. Normalise $\mathbf{v}=(1,1,1)$:
        $\|\mathbf{v}\|=\sqrt3$, $\mathbf{u}=(1,1,1)/\sqrt3$. Then
        $D_{\mathbf{u}}f=(1+1+1)/\sqrt3=\sqrt3$.

  \item $f_x=(\ln y)/x$, $f_y=(\ln x)/y$.

  \item $F_K=1+L^{1/2}/K^{1/2}$, $F_L=1+K^{1/2}/L^{1/2}$. From
        $F=(K^{1/2}+L^{1/2})^2$ with $g=K^{1/2}+L^{1/2}$:
        $F_K=g/K^{1/2}=(K^{1/2}+L^{1/2})/K^{1/2}=1+L^{1/2}/K^{1/2}$.
        \checkmark Same for $F_L$.

  \item $f_x=2xy$, $f_y=x^2+3y^2$. At $(-1,2)$: $\nabla f=(-4,1+12)^T
        =(-4,13)^T$. Normalise: $\|\mathbf{v}\|=5$,
        $\mathbf{u}=(4/5,-3/5)^T$. Then
        \[
          D_{\mathbf{u}}f=-4\cdot\frac45+13\cdot\left(-\frac35\right)
            =-\frac{16+39}{5}=-11 .
        \]

  \item $U_x=\alpha/x$, $U_y=\beta/y$.
        $\mathrm{MRS}_{xy}=\dfrac{\alpha/x}{\beta/y}=\dfrac{\alpha y}{\beta x}$.
        At $(\alpha,\beta)$: $\mathrm{MRS}=(\alpha\cdot\beta)/(\beta\cdot\alpha)=1$.

  \item $f=(x^2+y^2)^{-1/2}$. Chain rule:
        \[
          f_x=-\frac12(x^2+y^2)^{-3/2}\cdot2x=\frac{-x}{(x^2+y^2)^{3/2}},
          \qquad f_y=\frac{-y}{(x^2+y^2)^{3/2}} .
        \]
        So $\nabla f=-(x,y)^T/(x^2+y^2)^{3/2}$.

  \item Chain rule with $u=xy$:
        \[
          f_x=\frac{y}{1+x^2y^2}, \qquad f_y=\frac{x}{1+x^2y^2} .
        \]

  \item Let $S=K^\rho+L^\rho$, $F=S^{1/\rho}$. Then
        $F_K=\frac1\rho S^{1/\rho-1}\cdot\rho K^{\rho-1}
        =K^{\rho-1}S^{1/\rho-1}$. Since $S^{1/\rho}=F$, we have $S=F^\rho$
        and so $S^{1/\rho-1}=F^{1-\rho}\cdot F^{-1}\cdot F=F^{1-\rho}$.
        Thus $F_K=K^{\rho-1}F^{1-\rho}=(F/K)^{1-\rho}=(K/F)^{\rho-1}$.
        \checkmark

  \item $f_x=yz$, $f_y=xz$, $f_z=xy$. At $(2,3,-1)$:
        $\nabla f=(3\cdot(-1),2\cdot(-1),2\cdot3)^T=(-3,-2,6)^T$.
        Normalise $\mathbf{v}=(2,-1,2)$: $\|\mathbf{v}\|=3$,
        $\mathbf{u}=(2/3,-1/3,2/3)^T$. Then
        \[
          D_{\mathbf{u}}f=-3\cdot\frac23+(-2)\cdot\left(-\frac13\right)
            +6\cdot\frac23=-2+\frac23+4=2+\frac23=\frac83 .
        \]
\end{solutions}

\section{Additional Problems}

\begin{problems}
  \item Compute $f_x$ and $f_y$ for $f(x,y)=3x^5-2x^3y^2+4y^4$.
  \item For $F(K,L)=20\,K^{1/3}L^{2/3}$, find $F_K$ and $F_L$ at
        $(K,L)=(8,27)$.
  \item Compute $f_x$ and $f_y$ for $f(x,y)=e^{xy}$.
  \item Compute $\nabla f$ for $f(x,y)=2x^2y-3xy^3+5$.
  \item Compute the directional derivative of $f(x,y)=x^2+2xy-y^2$ at
        $(1,2)$ in the direction $\mathbf{v}=(3,4)$.
  \item For $F(K,L)=8K^{3/4}L^{1/4}$, compute $\mathrm{MP}_K$ and
        $\mathrm{MP}_L$ at $(K,L)=(16,16)$.
  \item Compute $f_x$ and $f_y$ for $f(x,y)=\ln(2x+3y)$.
  \item For $U(x,y)=3\ln(x)+4\ln(y)$, derive $\mathrm{MRS}_{xy}$ and
        evaluate at $(2,3)$.
  \item Compute $f_x$ and $f_y$ for $f(x,y)=(x+y^2)^4$.
  \item Compute $\nabla f$ for $f(x,y,z)=xy+yz+zx$.
  \item For a CES production $F(K,L)=5(K^{1/3}+L^{1/3})^3$, find $F_K$.
  \item Compute the directional derivative of $f(x,y)=x^2y$ at $(2,1)$
        in the direction $\mathbf{v}=(1,-1)$.
  \item For $U(x,y)=x^\alpha y^\beta$ (general Cobb--Douglas,
        $\alpha,\beta>0$), derive $\mathrm{MRS}_{xy}$ and show it depends
        only on the ratio $y/x$.
  \item Compute $f_x$ and $f_y$ for $f(x,y)=\dfrac{x-y}{x+y}$.
  \item For $f(K,L)=K^{1/2}L^{1/2}$ at $(K,L)=(16,9)$, find the unit
        direction of steepest ascent and the maximum rate of increase.
  \item Classify the open positive orthant
        $Q=\{(x,y)\in\R^2 : x>0,\ y>0\}$: is it open, closed, or
        neither? Is it bounded? Is it compact?
  \item Classify the budget set
        $B=\{(x,y)\in\R^2 : 2x+3y\le12,\ x\ge0,\ y\ge0\}$: open,
        closed, or neither? Bounded? Compact?
  \item Classify $M=\{(x,y)\in\R^2 : xy\ge1,\ x>0\}$: open, closed, or
        neither? Bounded? Compact? \emph{Hint:} check whether a
        sequence in $M$ can converge to a point with $x=0$.
  \item Classify the open disc
        $D=\{(x,y)\in\R^2 : x^2+y^2<4\}$: open, closed, or neither?
        Bounded? Compact?
  \item Classify (a) the half-open strip
        $S=\{(x,y)\in\R^2 : 0\le x<1\}$ and (b) the finite set
        $E=\{(0,0),\ (1,2),\ (3,1)\}$: open, closed, or neither?
        Bounded? Compact?
  \item Does the Weierstrass theorem guarantee a maximiser for the
        continuous function $u(x,y)=x^{1/2}y^{1/2}$ (a) on the budget
        set $B=\{(x,y): 2x+3y\le12,\ x\ge0,\ y\ge0\}$, (b) on the open
        positive orthant $Q=\{(x,y): x>0,\ y>0\}$? Justify each
        answer.
\end{problems}

In problems 22--29, for each function: state the domain; compute the
partials and the set on which they exist; decide, with a justification,
where $f$ is differentiable; and describe the graph near the points where
differentiability fails.

\begin{problems}[resume]
  \item $f(x,y)=\sqrt{(x-1)^{2}+(y+2)^{2}}$.
  \item $f(x,y)=\sqrt{4x^{2}+(y-3)^{2}}$.
  \item $f(x,y)=|x-2|+|y|$.
  \item $f(x,y)=\min\{2x,\,5y\}$ on $x,y>0$.
  \item $f(x,y)=\dfrac{3xy}{2x^{2}+y^{2}}$ for $(x,y)\neq(0,0)$, with
        $f(0,0)=0$. Also decide whether $f$ is continuous at the origin.
  \item $f(x,y)=\sqrt{x^{4}+y^{4}}$. The formula for $f_{x}$ breaks down
        at the origin; decide what really happens there.
  \item $f(x,y)=\sqrt{9-x^{2}-y^{2}}$.
  \item $f(x,y)=|x-y|$.
\end{problems}

In problems 30--37, find the tangent plane of $f$ at $P$ (the tangent
hyperplane, for three variables), then evaluate both $f$ and the plane at
the nearby point $Q$ and compare. Give both values to six decimals.

\begin{problems}[resume]
  \item $f(x,y)=\ln(2x+y^{2})$, \;$P=(0,1)$, \;$Q=(0.05,\,0.98)$.
  \item $f(x,y)=\mathrm{e}^{x-2y}$, \;$P=(2,1)$, \;$Q=(2.1,\,1.02)$.
  \item $f(x,y)=xy^{2}\,\mathrm{e}^{x-1}$, \;$P=(1,2)$,
        \;$Q=(0.98,\,2.03)$.
  \item $f(x,y)=\dfrac{x+y}{x-y}$, \;$P=(3,1)$, \;$Q=(3.1,\,0.9)$.
  \item $F(K,L)=K^{1/4}L^{3/4}$, \;$P=(16,1)$, \;$Q=(16.5,\,1.1)$; read
        the plane as an output estimate.
  \item $f(x,y,z)=\ln(x+2y+3z)$, \;$P=(2,-1,\tfrac13)$,
        \;$Q=(2.02,\,-0.99,\,0.35)$.
  \item $f(x,y,z)=x\,\mathrm{e}^{yz}$, \;$P=(2,0,5)$,
        \;$Q=(2.05,\,0.01,\,4.9)$.
  \item $f(x,y,z)=\dfrac{xy}{z}$, \;$P=(2,3,6)$, \;$Q=(2.1,\,2.9,\,6.2)$.
\end{problems}

In problems 38--46, compute the partial derivatives, and give the gradient
(a column) and the derivative (a row) at the stated point.

\begin{problems}[resume]
  \item $f(x_{1},\ldots,x_{n})=\sum_{i=1}^{n}\ln(1+x_{i}^{2})$, at
        $(1,1,\ldots,1)$.
  \item $f(x_{1},\ldots,x_{n})=\Bigl(\sum_{i=1}^{n}x_{i}\Bigr)^{3}$; then
        for $n=4$ at $(1,2,3,4)$.
  \item $f(x_{1},\ldots,x_{n})=\sum_{i=1}^{n}x_{i}^{\,i+1}$, at
        $(1,1,\ldots,1)$; then for $n=3$ at $(2,1,-1)$.
  \item $f(x_{1},\ldots,x_{n})=\sum_{j=1}^{n}\dfrac{j^{2}}{x_{j}^{3}}$ with
        all $x_{j}>0$, at $(1,1,\ldots,1)$.
  \item $f(x_{1},\ldots,x_{n})=\sum_{i=1}^{n}x_{i}^{2}
        -\dfrac1n\Bigl(\sum_{i=1}^{n}x_{i}\Bigr)^{2}$; show that the
        gradient vanishes exactly when all $x_{i}$ are equal, then evaluate
        for $n=3$ at $(1,2,6)$.
  \item $f(x,y,z)=y^{2}\ln(x^{2}+3yz)$, at $P=(1,1,1)$.
  \item $f(x,y,z)=x^{3}\,\mathrm{e}^{y^{2}-2xz}$, at $P=(1,0,1)$.
  \item $f(x,y)=2^{x}3^{y}+5^{xy}$, at $(1,1)$.
  \item $f(x_{1},\ldots,x_{n})=\sum_{i=1}^{n}3^{\,i x_{i}}$, at
        $(0,0,\ldots,0)$.
\end{problems}

In problems 47--52, decide differentiability directly from the definition,
by forming the remainder and estimating it, as in the third worked example.

\begin{problems}[resume]
  \item $f(x,y)=3xy$, at every point $(a,b)$.
  \item $f(x,y)=x^{2}+y^{2}$, first at the origin, then at every point.
  \item $f(x,y)=xy^{2}$, at every point $(a,b)$.
  \item $f(x,y)=y\,\mathrm{e}^{3x}$, at every point $(a,b)$.
  \item At the origin: (a) $f(x,y)=(x^{2}+4y^{2})^{1/3}$;
        (b) $g(x,y)=(x^{2}+y^{2})^{3/4}$.
  \item $f(x,y)=|xy|$: at the origin, and at the point $(2,0)$.
\end{problems}

In problems 53--54, $f(0,0)=0$. Show with the difference quotient that both
partials exist at the origin and are $0$; compute the limit along every
straight line through the origin and along the stated parabola; and decide
continuity and differentiability at the origin.

\begin{problems}[resume]
  \item $f(x,y)=\dfrac{x^{2}y}{x^{4}+y^{2}}$ for $(x,y)\neq(0,0)$; parabola
        $y=x^{2}$.
  \item $f(x,y)=\dfrac{2xy^{2}}{x^{2}+y^{4}}$ for $(x,y)\neq(0,0)$;
        parabola $x=y^{2}$.
\end{problems}

Problems 55--56 refer to the three surfaces below, each drawn over a square
around its special point.

\begin{gallery}
\begin{tikzpicture}
  \begin{axis}[surfaxis, title={(a)}, z post scale=1.1]
    \addplot3[surf, domain=-1:3, y domain=-2:2]{(x-1)^2+2*y^2};
  \end{axis}
\end{tikzpicture}\hfill
\begin{tikzpicture}
  \begin{axis}[surfaxis, title={(b)}, z post scale=1.1]
    \addplot3[surf, domain=-1:3, y domain=-2:2]{sqrt((x-1)^2+4*y^2)};
  \end{axis}
\end{tikzpicture}\hfill
\begin{tikzpicture}
  \begin{axis}[surfaxis, title={(c)}, z post scale=1.1]
    \addplot3[surf, domain=-2:2, y domain=-2:2]{2-abs(x)-abs(y)};
  \end{axis}
\end{tikzpicture}
\end{gallery}

\begin{problems}[resume]
  \item For each surface (a)--(c), say from the picture alone where the
        function is differentiable and where it is not, and why.
  \item Surface (b) has its tip above the point $(1,0)$, and its contours
        are ellipses twice as long in $x$ as in $y$. Which of the formulas
        $(x-1)^{2}+4y^{2}$, $\sqrt{(x-1)^{2}+4y^{2}}$ and
        $\sqrt{4(x-1)^{2}+y^{2}}$ can it be the graph of? Confirm the
        verdict of the previous problem for that formula with the
        difference quotient at $(1,0)$.
\end{problems}

\section{Additional Solutions}

\begin{solutions}
  \item $f_x=15x^4-6x^2y^2$, $f_y=-4x^3y+16y^3$.

  \item $F_K=20\cdot\frac13K^{-2/3}L^{2/3}=\frac{20}{3}(L/K)^{2/3}$;
        $F_L=20\cdot\frac23K^{1/3}L^{-1/3}=\frac{40}{3}(K/L)^{1/3}$. At
        $(8,27)$: $(L/K)^{2/3}=(27/8)^{2/3}=9/4$, so $F_K=15$;
        $(K/L)^{1/3}=(8/27)^{1/3}=2/3$, so
        $F_L=\frac{40}{3}\cdot\frac23=\frac{80}{9}$.

  \item $f_x=ye^{xy}$, $f_y=xe^{xy}$.

  \item $\nabla f=(4xy-3y^3,\ 2x^2-9xy^2)^T$.

  \item $f_x=2x+2y$, $f_y=2x-2y$. At $(1,2)$: $\nabla f=(6,-2)^T$.
        Normalise: $\|\mathbf{v}\|=5$, $\mathbf{u}=(3/5,4/5)^T$. Then
        \[
          D_{\mathbf{u}}f=6\cdot\frac35+(-2)\cdot\frac45
            =\frac{18-8}{5}=2 .
        \]

  \item $F_K=8\cdot\frac34K^{-1/4}L^{1/4}=6(L/K)^{1/4}$,
        $F_L=2(K/L)^{3/4}$. At $(16,16)$: $\mathrm{MP}_K=6$,
        $\mathrm{MP}_L=2$.

  \item $f_x=\dfrac{2}{2x+3y}$, $f_y=\dfrac{3}{2x+3y}$.

  \item $U_x=3/x$, $U_y=4/y$.
        $\mathrm{MRS}_{xy}=U_x/U_y=\dfrac{3y}{4x}$. At $(2,3)$:
        $\mathrm{MRS}=\dfrac98$.

  \item Chain: $f_x=4(x+y^2)^3$,
        $f_y=4(x+y^2)^3\cdot2y=8y(x+y^2)^3$.

  \item $f_x=y+z$, $f_y=x+z$, $f_z=x+y$. So
        $\nabla f=(y+z,\ x+z,\ x+y)^T$.

  \item Let $g=K^{1/3}+L^{1/3}$, so $F=5g^3$. Then
        $F_K=5\cdot3g^2\cdot\frac13K^{-2/3}=5g^2K^{-2/3}
        =\dfrac{5(K^{1/3}+L^{1/3})^2}{K^{2/3}}$.

  \item $f_x=2xy$, $f_y=x^2$. At $(2,1)$: $\nabla f=(4,4)^T$. Normalise:
        $\|\mathbf{v}\|=\sqrt2$, $\mathbf{u}=(1/\sqrt2,-1/\sqrt2)^T$.
        Then
        \[
          D_{\mathbf{u}}f=4\cdot\frac{1}{\sqrt2}
            +4\cdot\left(-\frac{1}{\sqrt2}\right)=0 .
        \]
        (Tangent to the level curve.)

  \item $U_x=\alpha x^{\alpha-1}y^\beta$,
        $U_y=\beta x^\alpha y^{\beta-1}$. So
        \[
          \mathrm{MRS}_{xy}=\frac{U_x}{U_y}
            =\frac{\alpha}{\beta}\cdot\frac{y}{x} .
        \]
        Depends only on $y/x$ (and the exponent ratio), not on the
        scale. \checkmark

  \item Quotient rule with $D=x+y$:
        \[
          f_x=\frac{D-(x-y)}{D^2}=\frac{2y}{(x+y)^2}, \qquad
          f_y=\frac{-D-(x-y)}{D^2}=\frac{-2x}{(x+y)^2} .
        \]

  \item $f_K=\frac12K^{-1/2}L^{1/2}$, $f_L=\frac12K^{1/2}L^{-1/2}$. At
        $(16,9)$: $f_K=\frac12\cdot\frac34=3/8$;
        $f_L=\frac12\cdot\frac43=2/3$. So $\nabla f=(3/8,2/3)^T$.
        Magnitude:
        $\|\nabla f\|=\sqrt{9/64+4/9}=\sqrt{81/576+256/576}
        =\sqrt{337/576}=\sqrt{337}/24\approx0.765$. Unit direction of
        steepest ascent:
        $\mathbf{u}^*=(3/8,2/3)^T/(\sqrt{337}/24)=(9,16)^T/\sqrt{337}$.

  \item \textbf{Open.} For any $(x,y)\in Q$ the ball of radius
        $r=\min\{x,y\}>0$ around $(x,y)$ stays inside $Q$, so every
        point is interior. \textbf{Not closed:} the boundary points
        $(x,0)$ and $(0,y)$ with $x,y\ge0$ (the coordinate half-axes,
        including the origin) do not belong to $Q$. \textbf{Not
        bounded:} $Q$ contains $(t,t)$ for every $t>0$, and
        $\|(t,t)\|=t\sqrt2\to\infty$. Hence $Q$ is open, unbounded and
        \textbf{not compact}.

  \item \textbf{Closed.} $B$ is the intersection of the three sets
        $\{2x+3y\le12\}$, $\{x\ge0\}$ and $\{y\ge0\}$, each defined by
        a non-strict inequality on a continuous function; each is
        closed, and an intersection of closed sets is closed (every
        boundary point of $B$ satisfies all three inequalities, so it
        lies in $B$). \textbf{Not open:} $(0,0)\in B$, but every ball
        around the origin contains points with negative coordinates,
        which are outside $B$; so $(0,0)$ is not interior.
        \textbf{Bounded:} from the constraints, $0\le x\le6$ and
        $0\le y\le4$, so $B\subseteq B_8(\mathbf{0})$. Closed and
        bounded in $\R^2$: $B$ is \textbf{compact}.

  \item First note the constraint $x>0$ together with $xy\ge1$ forces
        $y\ge1/x>0$. \textbf{Closed:} suppose
        $(x_n,y_n)\in M$ with $(x_n,y_n)\to(a,b)$. If $a=0$, then
        $y_n\to b$ finite would give $x_ny_n\to0<1$, contradicting
        $x_ny_n\ge1$; so $a>0$, and passing to the limit in
        $x_ny_n\ge1$ gives $ab\ge1$, hence $(a,b)\in M$. $M$ contains
        all its limit (boundary) points, so it is closed: even
        though it is written with a strict inequality, that inequality
        is never active. \textbf{Not open:} $(1,1)\in M$ lies on the
        hyperbola $xy=1$, and every ball around it contains points
        with $xy<1$. \textbf{Not bounded:} $(t,t)\in M$ for all
        $t\ge1$. Hence $M$ is closed, unbounded and \textbf{not
        compact}.

  \item \textbf{Open.} For $(x,y)\in D$ put
        $r=2-\sqrt{x^2+y^2}>0$; the ball $B_r(x,y)$ lies in $D$ by the
        triangle inequality. \textbf{Not closed:} the boundary circle
        $\{x^2+y^2=4\}$ is not contained in $D$ (e.g.\ $(2,0)\notin
        D$). \textbf{Bounded:} $D\subseteq B_3(\mathbf{0})$. Bounded
        but not closed: $D$ is \textbf{not compact}.

  \item (a) \textbf{Neither open nor closed.} The points $(0,y)$
        belong to $S$ but are boundary points (every ball around them
        contains points with $x<0$), so $S$ is not open; the boundary
        points $(1,y)$ do not belong to $S$, so $S$ is not closed.
        \textbf{Unbounded:} $(0,t)\in S$ for all $t$. Hence
        \textbf{not compact}. (b) \textbf{Closed:} around any point
        not in $E$ there is a ball avoiding all three points (a
        finite set always contains all its boundary points:
        indeed $E$ is its own boundary). \textbf{Not open:} no ball
        around $(0,0)$ lies inside $E$. \textbf{Bounded:}
        $E\subseteq B_4(\mathbf{0})$. Closed and bounded: $E$ is
        \textbf{compact}.

  \item (a) \textbf{Yes.} $B$ is compact (closed and bounded, see the
        budget-set problem above) and $u$ is continuous on $B$, so by
        the Weierstrass theorem $u$ attains both a maximum and a
        minimum on $B$; a maximiser exists. (b) \textbf{No.} $Q$ is
        not compact (it is neither closed nor bounded), so Weierstrass
        does not apply: and indeed no maximiser exists:
        $u(t,t)=t\to\infty$ as $t\to\infty$, so $u$ is unbounded above
        on $Q$.

  \item \textbf{Domain:} the radicand is a sum of squares, never
        negative, so $D=\R^{2}$. \textbf{Partials:} where the radicand is
        positive, i.e.\ for $(x,y)\neq(1,-2)$, the chain rule gives
        \[
          f_{x}=\frac{x-1}{\sqrt{(x-1)^{2}+(y+2)^{2}}},\qquad
          f_{y}=\frac{y+2}{\sqrt{(x-1)^{2}+(y+2)^{2}}} .
        \]
        At $(1,-2)$ use the quotient:
        $\bigl(f(1+h,-2)-f(1,-2)\bigr)/h=\sqrt{h^{2}}/h=|h|/h$, which is $1$
        for $h>0$ and $-1$ for $h<0$, so it has no limit. $f_{x}(1,-2)$
        does not exist, and in the same way neither does $f_{y}(1,-2)$. The
        partials exist exactly on $\R^{2}\setminus\{(1,-2)\}$.
        \textbf{Differentiability:} on that open set the partials are
        continuous ($f$ is $C^{1}$), so $f$ is differentiable there. At
        $(1,-2)$ it is not, since differentiability would force the
        partials to exist. Note that $f$ \emph{is} continuous at $(1,-2)$:
        continuity is not enough. \textbf{Graph:} a circular cone with its
        vertex at $(1,-2,0)$, opening upwards with slope $1$ in every
        direction; the contours are circles of radius $c$ about $(1,-2)$.
        At the sharp tip there is no tangent plane.

  \item \textbf{Domain:} $\R^{2}$, the radicand being a sum of squares.
        \textbf{Partials:} for $(x,y)\neq(0,3)$,
        \[
          f_{x}=\frac{4x}{\sqrt{4x^{2}+(y-3)^{2}}},\qquad
          f_{y}=\frac{y-3}{\sqrt{4x^{2}+(y-3)^{2}}} .
        \]
        At $(0,3)$: $f(h,3)-f(0,3)=\sqrt{4h^{2}}=2|h|$, so the quotient
        $2|h|/h=\pm2$ has no limit; and $f(0,3+h)-f(0,3)=|h|$ gives
        $\pm1$. Neither partial exists at $(0,3)$; both exist everywhere
        else. \textbf{Differentiability:} $C^{1}$, hence differentiable, on
        $\R^{2}\setminus\{(0,3)\}$; not differentiable at $(0,3)$.
        \textbf{Graph:} an elliptic cone with apex $(0,3,0)$. The contours
        $4x^{2}+(y-3)^{2}=c^{2}$ are ellipses about $(0,3)$ with semi-axes
        $c/2$ in the $x$-direction and $c$ in the $y$-direction: slope $2$
        across $x$, slope $1$ along $y$, and a sharp point at the apex.

  \item \textbf{Domain:} $\R^{2}$. \textbf{Partials:} for $x\neq2$,
        $f_{x}=1$ if $x>2$ and $f_{x}=-1$ if $x<2$; at a point $(2,y_{0})$,
        $\bigl(f(2+h,y_{0})-f(2,y_{0})\bigr)/h=|h|/h=\pm1$, so $f_{x}$ does
        not exist anywhere on the line $x=2$. Likewise $f_{y}=\pm1$ (the
        sign of $y$) for $y\neq0$, and $f_{y}$ does not exist on the line
        $y=0$. Both partials exist exactly off the two lines $x=2$ and
        $y=0$. \textbf{Differentiability:} off the two lines $f$ is, near
        each point, one of the four linear functions $\pm(x-2)\pm y$, so it
        is differentiable there; on the lines it is not, since a partial
        is missing. \textbf{Graph:} an upside-down square pyramid with its
        vertex at $(2,0,0)$; the contours $|x-2|+|y|=c$ are diamonds with
        corners $(2\pm c,0)$ and $(2,\pm c)$, and the surface is creased
        above the lines $x=2$ and $y=0$.

  \item \textbf{Domain:} the open quadrant $x,y>0$, as given.
        \textbf{Partials:} where $2x<5y$ (above the ray
        $y=\tfrac25x$), $f=2x$, so $f_{x}=2$, $f_{y}=0$. Where $2x>5y$,
        $f=5y$, so $f_{x}=0$, $f_{y}=5$. On the ray, at $(x_{0},y_{0})$ with
        $2x_{0}=5y_{0}$:
        \[
          f(x_{0}+h,y_{0})-f(x_{0},y_{0})=\min\{2x_{0}+2h,\,5y_{0}\}-5y_{0}
          =\min\{2h,\,0\},
        \]
        so the quotient is $0$ for $h>0$ and $2$ for $h<0$: no limit, no
        $f_{x}$. Similarly $f(x_{0},y_{0}+k)-f(x_{0},y_{0})=\min\{0,5k\}$
        gives one-sided slopes $0$ and $5$: no $f_{y}$. The partials exist
        exactly off the ray $y=\tfrac25x$. \textbf{Differentiability:}
        differentiable off the ray (locally linear), not on it.
        \textbf{Graph:} two planes, $z=2x$ and $z=5y$, meeting in a ridge
        above the ray: a roof. The contours are L-shaped corners, the level
        $c$ having its corner at $(c/2,\,c/5)$ on the ray, which is the
        Leontief isoquant picture.

  \item \textbf{Domain:} the denominator $2x^{2}+y^{2}$ vanishes only at
        the origin, where $f$ is defined separately, so $D=\R^{2}$.
        \textbf{Partials off the origin}, quotient rule with
        $N=2x^{2}+y^{2}$:
        \[
          f_{x}=\frac{3y\,N-3xy\cdot4x}{N^{2}}
               =\frac{3y\,(y^{2}-2x^{2})}{(2x^{2}+y^{2})^{2}},\qquad
          f_{y}=\frac{3x\,N-3xy\cdot2y}{N^{2}}
               =\frac{3x\,(2x^{2}-y^{2})}{(2x^{2}+y^{2})^{2}} .
        \]
        \textbf{At the origin}, from the definition: $f(h,0)=0$ and
        $f(0,h)=0$ for all $h$, so both quotients are $0$ and
        $f_{x}(0,0)=f_{y}(0,0)=0$. The partials exist \emph{everywhere}.
        \textbf{Continuity:} on the line $y=mx$, $x\neq0$,
        \[
          f(x,mx)=\frac{3mx^{2}}{(2+m^{2})x^{2}}=\frac{3m}{2+m^{2}},
        \]
        a constant that depends on $m$; on $y=x$ it is $1\neq f(0,0)$. So
        $f$ is not continuous at the origin, and hence not differentiable
        there, even though both partials exist. Away from the origin it is
        $C^{1}$ and differentiable. \textbf{Graph:} each ray from the origin
        carries its own constant height $3m/(2+m^{2})$, ranging between
        $-\tfrac{3\sqrt2}{4}$ and $\tfrac{3\sqrt2}{4}$ (attained at
        $m=\mp\sqrt2$); all these heights meet above the origin, so the
        surface is a twisted sheet torn at that point.

  \item \textbf{Domain:} $\R^{2}$. \textbf{Partials} off the origin:
        $f_{x}=\dfrac{2x^{3}}{\sqrt{x^{4}+y^{4}}}$,
        $f_{y}=\dfrac{2y^{3}}{\sqrt{x^{4}+y^{4}}}$. At the origin the
        formula fails, so use the quotient:
        $f(h,0)=\sqrt{h^{4}}=h^{2}$, and $h^{2}/h=h\to0$. So
        $f_{x}(0,0)=0$, and likewise $f_{y}(0,0)=0$: the partials exist
        everywhere, despite the broken formula. \textbf{Differentiability
        at the origin:} with $\mathbf{a}=\mathbf{0}$ the remainder is
        $r(\mathbf{h})=\sqrt{h_{1}^{4}+h_{2}^{4}}$. Since
        $h_{1}^{4}+h_{2}^{4}\le(h_{1}^{2}+h_{2}^{2})^{2}$ (the missing
        cross term $2h_{1}^{2}h_{2}^{2}$ is non-negative),
        \[
          0\le\frac{r(\mathbf{h})}{\|\mathbf{h}\|}
          \le\frac{h_{1}^{2}+h_{2}^{2}}{\|\mathbf{h}\|}=\|\mathbf{h}\|\to0 .
        \]
        So $f$ is differentiable at the origin with tangent plane $z=0$,
        and, being $C^{1}$ elsewhere, differentiable on all of $\R^{2}$.
        \textbf{Graph:} a bowl with a flat bottom at the origin; its
        contours $x^{4}+y^{4}=c^{2}$ are rounded squares. Nothing fails.

  \item \textbf{Domain:} $9-x^{2}-y^{2}\ge0$, the closed disc of radius $3$.
        \textbf{Partials:} for $x^{2}+y^{2}<9$,
        \[
          f_{x}=\frac{-x}{\sqrt{9-x^{2}-y^{2}}},\qquad
          f_{y}=\frac{-y}{\sqrt{9-x^{2}-y^{2}}},
        \]
        continuous on the open disc. On the boundary circle the radicand
        is $0$, the formulas divide by $0$, and, more basically, the points
        are not interior: every neighbourhood of a boundary point contains
        points outside the disc, where $f$ is not defined, so the two-sided
        difference quotients needed for the partials do not exist. \textbf{Differentiability:}
        $f$ is differentiable exactly on the open disc $x^{2}+y^{2}<9$.
        \textbf{Graph:} the upper hemisphere of radius $3$. Its tangent
        plane is horizontal at the top $(0,0,3)$, tilts more and more
        steeply outwards, and becomes vertical at the rim, where
        $\|\grad f\|\to\infty$.

  \item \textbf{Domain:} $\R^{2}$. \textbf{Partials:} where $x>y$,
        $f=x-y$, so $f_{x}=1$, $f_{y}=-1$; where $x<y$, $f=y-x$, so
        $f_{x}=-1$, $f_{y}=1$. At a point $(x_{0},x_{0})$ of the diagonal,
        $f(x_{0}+h,x_{0})-f(x_{0},x_{0})=|h|$, and $|h|/h=\pm1$ has no limit;
        the same holds in $y$. The partials exist exactly off the line
        $y=x$. \textbf{Differentiability:} differentiable off $y=x$, not on
        it, although $f$ is continuous everywhere. \textbf{Graph:} a
        V-shaped valley whose crease, at height $0$, runs along the line
        $y=x$; the contours are pairs of parallel lines $y=x\pm c$, the
        tilted-trough picture of Chapter~9 with a sharp floor.

  \item $f(P)=\ln(0+1)=0$. Chain rule:
        $f_{x}=\dfrac{2}{2x+y^{2}}$, $f_{y}=\dfrac{2y}{2x+y^{2}}$, so at
        $P$ (denominator $1$): $f_{x}=2$, $f_{y}=2$. Tangent plane:
        \[
          z=0+2(x-0)+2(y-1)=2x+2y-2 .
        \]
        At $Q$: $f(Q)=\ln(0.1+0.9604)=\ln1.0604\approx0.058646$, and
        $L(Q)=0.1+1.96-2=0.060000$. The plane overestimates by about
        $0.001354$.

  \item $f(P)=\mathrm{e}^{0}=1$. $f_{x}=\mathrm{e}^{x-2y}$,
        $f_{y}=-2\mathrm{e}^{x-2y}$, so at $P$: $f_{x}=1$, $f_{y}=-2$.
        Tangent plane:
        \[
          z=1+(x-2)-2(y-1)=x-2y+1 .
        \]
        At $Q$: $f(Q)=\mathrm{e}^{2.1-2.04}=\mathrm{e}^{0.06}\approx1.061837$,
        and $L(Q)=2.1-2.04+1=1.060000$. The plane underestimates by about
        $0.001837$: the exponential is convex, so it lies above all its
        tangent planes.

  \item $f(P)=1\cdot4\cdot\mathrm{e}^{0}=4$. Product rule in $x$, where
        both $x$ and $\mathrm{e}^{x-1}$ depend on $x$:
        \[
          f_{x}=y^{2}\mathrm{e}^{x-1}+xy^{2}\mathrm{e}^{x-1}
               =y^{2}\mathrm{e}^{x-1}(1+x),\qquad
          f_{y}=2xy\,\mathrm{e}^{x-1} .
        \]
        At $P$: $f_{x}=4\cdot1\cdot2=8$, $f_{y}=2\cdot1\cdot2\cdot1=4$.
        Tangent plane:
        \[
          z=4+8(x-1)+4(y-2) .
        \]
        At $Q$: $f(Q)=0.98\cdot2.03^{2}\cdot\mathrm{e}^{-0.02}
        =0.98\cdot4.1209\cdot0.980199\approx3.958515$, and
        $L(Q)=4-0.16+0.12=3.960000$. The plane is off by about
        $0.001485$.

  \item $f(P)=\dfrac{4}{2}=2$. Quotient rule with denominator $(x-y)^{2}$:
        \[
          f_{x}=\frac{(x-y)-(x+y)}{(x-y)^{2}}=\frac{-2y}{(x-y)^{2}},\qquad
          f_{y}=\frac{(x-y)+(x+y)}{(x-y)^{2}}=\frac{2x}{(x-y)^{2}} .
        \]
        At $P$: $f_{x}=-\tfrac24=-\tfrac12$, $f_{y}=\tfrac64=\tfrac32$.
        Tangent plane:
        \[
          z=2-\tfrac12(x-3)+\tfrac32(y-1) .
        \]
        At $Q$: $f(Q)=\dfrac{4.0}{2.2}\approx1.818182$, and
        $L(Q)=2-0.05-0.15=1.800000$. The error, about $0.018182$, is
        several times larger than in the neighbouring problems. The step
        keeps $x+y=4$ but moves the denominator $x-y$ from $2$ to $2.2$,
        and along that direction $f=4/(x-y)$ is strongly curved. The plane
        is a local approximation, and its quality depends on the curvature
        as much as on the length of the step.

  \item $F(P)=16^{1/4}\cdot1=2$.
        $F_{K}=\tfrac14K^{-3/4}L^{3/4}$, $F_{L}=\tfrac34K^{1/4}L^{-1/4}$; at
        $P$: $F_{K}=\tfrac14\cdot\tfrac18=\tfrac1{32}$ and
        $F_{L}=\tfrac34\cdot2=\tfrac32$. Tangent plane:
        \[
          z=2+\tfrac1{32}(K-16)+\tfrac32(L-1) .
        \]
        Economically: output $2$ plus marginal product times input change,
        for each input. At $Q$: $L(Q)=2+\tfrac{0.5}{32}+\tfrac32\cdot0.1
        =2+0.015625+0.15=2.165625$, while
        $F(Q)=16.5^{1/4}\cdot1.1^{3/4}\approx2.164789$. The estimate is
        high by about $0.000836$: with exponents summing to $1$ this
        Cobb--Douglas function is concave, so its graph lies below every
        tangent plane.

  \item $f(P)=\ln(2-2+1)=\ln1=0$. With $s=x+2y+3z$:
        $f_{x}=\tfrac1s$, $f_{y}=\tfrac2s$, $f_{z}=\tfrac3s$, and $s=1$ at
        $P$, so $\grad f(P)=(1,2,3)^{T}$. Tangent hyperplane:
        \[
          w=0+(x-2)+2(y+1)+3\bigl(z-\tfrac13\bigr)=x+2y+3z-1 .
        \]
        At $Q$: $s=2.02-1.98+1.05=1.09$, so $f(Q)=\ln1.09\approx0.086178$,
        and $L(Q)=1.09-1=0.090000$. The hyperplane overestimates by about
        $0.003822$, since $\ln$ is concave.

  \item $f(P)=2\mathrm{e}^{0}=2$. $f_{x}=\mathrm{e}^{yz}$,
        $f_{y}=xz\,\mathrm{e}^{yz}$, $f_{z}=xy\,\mathrm{e}^{yz}$; at $P$:
        $f_{x}=1$, $f_{y}=2\cdot5=10$, $f_{z}=2\cdot0=0$. Tangent
        hyperplane:
        \[
          w=2+(x-2)+10(y-0)+0\cdot(z-5)=x+10y .
        \]
        The change in $z$ does not enter at all, because $f_{z}(P)=0$. At
        $Q$: $f(Q)=2.05\,\mathrm{e}^{0.049}\approx2.152952$, and
        $L(Q)=2.05+0.1=2.150000$; error about $0.002952$.

  \item $f(P)=\tfrac66=1$. $f_{x}=\tfrac yz$, $f_{y}=\tfrac xz$,
        $f_{z}=-\tfrac{xy}{z^{2}}$; at $P$: $f_{x}=\tfrac12$,
        $f_{y}=\tfrac13$, $f_{z}=-\tfrac6{36}=-\tfrac16$. Tangent
        hyperplane:
        \[
          w=1+\tfrac12(x-2)+\tfrac13(y-3)-\tfrac16(z-6) .
        \]
        At $Q$: $L(Q)=1+0.05-0.033333-0.033333=0.983333$, and
        $f(Q)=\dfrac{2.1\cdot2.9}{6.2}=\dfrac{6.09}{6.2}\approx0.982258$;
        error about $0.001075$.

  \item Only the term $i=k$ contains $x_{k}$, so
        \[
          \pd{f}{x_{k}}=\frac{2x_{k}}{1+x_{k}^{2}},\qquad k=1,\ldots,n .
        \]
        At $(1,\ldots,1)$ every partial is $\tfrac22=1$:
        $\grad f=(1,1,\ldots,1)^{T}$ and $Df=(1\ \ 1\ \ \cdots\ \ 1)$.

  \item Chain rule with inner function $s=\sum_{i}x_{i}$, whose partial in
        $x_{k}$ is $1$:
        \[
          \pd{f}{x_{k}}=3\Bigl(\sum_{i=1}^{n}x_{i}\Bigr)^{2}\cdot1 ,
        \]
        the same for every $k$. For $n=4$ at $(1,2,3,4)$: $s=10$, so each
        partial is $300$, $\grad f=(300,300,300,300)^{T}$ and
        $Df=(300\ \ 300\ \ 300\ \ 300)$.

  \item The term containing $x_{k}$ is $x_{k}^{k+1}$, so
        $\pd{f}{x_{k}}=(k+1)\,x_{k}^{k}$. At $(1,\ldots,1)$:
        $\grad f=(2,3,\ldots,n+1)^{T}$, $Df=(2\ \ 3\ \ \cdots\ \ n+1)$. For
        $n=3$, $f=x_{1}^{2}+x_{2}^{3}+x_{3}^{4}$ with partials
        $2x_{1},\,3x_{2}^{2},\,4x_{3}^{3}$; at $(2,1,-1)$:
        $\grad f=(4,3,-4)^{T}$ and $Df=(4\ \ 3\ \ {-4})$.

  \item Write $f=\sum_{j}j^{2}x_{j}^{-3}$. The term with $j=k$ gives
        \[
          \pd{f}{x_{k}}=k^{2}\cdot(-3)x_{k}^{-4}=-\frac{3k^{2}}{x_{k}^{4}} .
        \]
        At $(1,\ldots,1)$: $\grad f=-3\,(1,4,9,\ldots,n^{2})^{T}$ and
        $Df=(-3\ \ {-12}\ \ {-27}\ \ \cdots\ \ {-3n^{2}})$.

  \item The first sum contributes $2x_{k}$; the second term, by the chain
        rule as in problem 39, contributes
        $-\tfrac1n\cdot2\sum_{i}x_{i}$. With the mean
        $\bar x=\tfrac1n\sum_{i}x_{i}$:
        \[
          \pd{f}{x_{k}}=2x_{k}-\frac2n\sum_{i=1}^{n}x_{i}=2\,(x_{k}-\bar x) .
        \]
        The gradient is zero exactly when $x_{k}=\bar x$ for every $k$,
        i.e.\ when all $x_{i}$ are equal. For $n=3$ at $(1,2,6)$:
        $\bar x=3$, so $\grad f=(-4,-2,6)^{T}$ and $Df=(-4\ \ {-2}\ \ 6)$.
        (The components always add up to $0$.)

  \item With $s=x^{2}+3yz$: $f_{x}=y^{2}\cdot\dfrac{2x}{s}$, since
        $y^{2}$ is a constant factor; $f_{y}$ needs the product rule,
        because $y$ appears in both factors:
        \[
          f_{y}=2y\ln s+y^{2}\cdot\frac{3z}{s};
        \]
        and $f_{z}=y^{2}\cdot\dfrac{3y}{s}=\dfrac{3y^{3}}{s}$. At
        $P=(1,1,1)$, $s=4$:
        \[
          f_{x}=\tfrac24=\tfrac12,\qquad
          f_{y}=2\ln4+\tfrac34\approx3.522589,\qquad
          f_{z}=\tfrac34 .
        \]
        So $\grad f(P)=\bigl(\tfrac12,\ 2\ln4+\tfrac34,\ \tfrac34\bigr)^{T}$
        and $Df(P)=\bigl(\tfrac12\ \ \ 2\ln4+\tfrac34\ \ \ \tfrac34\bigr)$.

  \item Write $E=\mathrm{e}^{y^{2}-2xz}$. In $x$ both factors move, so
        product and chain rule:
        \[
          f_{x}=3x^{2}E+x^{3}E\cdot(-2z)=x^{2}(3-2xz)\,E .
        \]
        In $y$ and $z$ only the exponential moves:
        $f_{y}=x^{3}E\cdot2y=2x^{3}y\,E$ and
        $f_{z}=x^{3}E\cdot(-2x)=-2x^{4}E$. At $P=(1,0,1)$,
        $E=\mathrm{e}^{-2}$:
        $f_{x}=1\cdot(3-2)\mathrm{e}^{-2}=\mathrm{e}^{-2}$, $f_{y}=0$,
        $f_{z}=-2\mathrm{e}^{-2}$. So
        $\grad f(P)=\mathrm{e}^{-2}(1,0,-2)^{T}\approx(0.135335,\,0,\,-0.270671)^{T}$
        and $Df(P)=\mathrm{e}^{-2}(1\ \ 0\ \ {-2})$.

  \item Use $\partial a^{g}/\partial x=a^{g}\ln a\cdot g_{x}$. In the first
        term $3^{y}$ is a constant factor; in the second the exponent $xy$
        has $x$-derivative $y$:
        \[
          f_{x}=2^{x}3^{y}\ln2+y\,5^{xy}\ln5,\qquad
          f_{y}=2^{x}3^{y}\ln3+x\,5^{xy}\ln5 .
        \]
        At $(1,1)$, $2^{x}3^{y}=6$ and $5^{xy}=5$:
        \[
          \grad f(1,1)=\begin{pmatrix}6\ln2+5\ln5\\6\ln3+5\ln5\end{pmatrix}
          \approx\begin{pmatrix}12.206073\\14.638863\end{pmatrix},
        \]
        and $Df(1,1)$ is the same pair written as a row.

  \item Only the term $i=k$ contains $x_{k}$, and its exponent $kx_{k}$ has
        derivative $k$:
        \[
          \pd{f}{x_{k}}=3^{\,kx_{k}}\ln3\cdot k=k\ln3\cdot3^{\,kx_{k}} .
        \]
        At the origin every $3^{0}=1$, so
        $\grad f=\ln3\,(1,2,\ldots,n)^{T}$ and
        $Df=\ln3\,(1\ \ 2\ \ \cdots\ \ n)$.

  \item Candidate: $f_{x}=3y$, $f_{y}=3x$, so $\grad f(a,b)=(3b,3a)$.
        Remainder:
        \[
          r(\mathbf{h})=3(a+h_{1})(b+h_{2})-3ab-3bh_{1}-3ah_{2}=3h_{1}h_{2} .
        \]
        Estimate:
        $\dfrac{|r(\mathbf{h})|}{\|\mathbf{h}\|}\le\dfrac{3\cdot\frac12
        \|\mathbf{h}\|^{2}}{\|\mathbf{h}\|}=\tfrac32\|\mathbf{h}\|\to0$.
        So $f$ is differentiable at every point, with derivative
        $(3b,\,3a)$.
  \item At the origin $\grad f=(0,0)$ and
        $r(\mathbf{h})=h_{1}^{2}+h_{2}^{2}=\|\mathbf{h}\|^{2}$, so
        $|r|/\|\mathbf{h}\|=\|\mathbf{h}\|\to0$: differentiable, with tangent
        plane $z=0$. At $(a,b)$: $\grad f=(2a,2b)$ and
        \[
          r(\mathbf{h})=(a+h_{1})^{2}+(b+h_{2})^{2}-a^{2}-b^{2}-2ah_{1}-2bh_{2}
          =h_{1}^{2}+h_{2}^{2},
        \]
        the same remainder, so the same estimate works everywhere.
  \item Candidate: $f_{x}=y^{2}$, $f_{y}=2xy$, so $\grad f(a,b)=(b^{2},2ab)$.
        Expand $(a+h_{1})(b^{2}+2bh_{2}+h_{2}^{2})
        =ab^{2}+2abh_{2}+ah_{2}^{2}+b^{2}h_{1}+2bh_{1}h_{2}+h_{1}h_{2}^{2}$
        and subtract $ab^{2}+b^{2}h_{1}+2abh_{2}$:
        \[
          r(\mathbf{h})=a\,h_{2}^{2}+2b\,h_{1}h_{2}+h_{1}h_{2}^{2} .
        \]
        Estimate:
        $\dfrac{|r|}{\|\mathbf{h}\|}\le|a|\,\|\mathbf{h}\|
        +|b|\,\|\mathbf{h}\|+\|\mathbf{h}\|^{2}\to0$. Differentiable
        everywhere, derivative $(b^{2},\,2ab)$.
  \item Candidate: $f_{x}=3y\mathrm{e}^{3x}$, $f_{y}=\mathrm{e}^{3x}$, so
        $\grad f(a,b)=\mathrm{e}^{3a}(3b,\,1)$. Write
        $\mathrm{e}^{3(a+h_{1})}=\mathrm{e}^{3a}\bigl(1+3h_{1}+R(3h_{1})\bigr)$:
        \[
          f(a+h_{1},b+h_{2})=\mathrm{e}^{3a}(b+h_{2})\bigl(1+3h_{1}+R\bigr)
          =\mathrm{e}^{3a}\bigl(b+3bh_{1}+h_{2}+3h_{1}h_{2}+(b+h_{2})R\bigr).
        \]
        Subtracting $b\mathrm{e}^{3a}$ and $\mathrm{e}^{3a}(3bh_{1}+h_{2})$
        leaves
        \[
          r(\mathbf{h})=\mathrm{e}^{3a}\bigl(3h_{1}h_{2}+(b+h_{2})\,R(3h_{1})\bigr).
        \]
        For $\|\mathbf{h}\|\le\tfrac13$ we have $|3h_{1}|\le1$, so
        $|R(3h_{1})|\le2\cdot9h_{1}^{2}\le18\|\mathbf{h}\|^{2}$, and
        \[
          \frac{|r|}{\|\mathbf{h}\|}\le\mathrm{e}^{3a}\Bigl(\tfrac32
          \|\mathbf{h}\|+18\bigl(|b|+\|\mathbf{h}\|\bigr)\|\mathbf{h}\|\Bigr)
          \to0 .
        \]
        Differentiable at every point.
  \item (a) On the $x$-axis $f(h,0)=(h^{2})^{1/3}=|h|^{2/3}$, so the
        difference quotient $|h|^{2/3}/h=\pm|h|^{-1/3}$ is unbounded:
        $f_{x}(0,0)$ does not exist (likewise $f_{y}$, with quotient
        $\pm4^{1/3}|h|^{-1/3}$). Not differentiable.
        (b) $g(\mathbf{h})=\|\mathbf{h}\|^{3/2}$. On an axis the quotient is
        $|h|^{3/2}/h=\pm|h|^{1/2}\to0$, so both partials are $0$. Then
        $r(\mathbf{h})=\|\mathbf{h}\|^{3/2}$ and
        $|r|/\|\mathbf{h}\|=\|\mathbf{h}\|^{1/2}\to0$: differentiable, with
        tangent plane $z=0$. The exponents decide: $2p=\tfrac23<1$ fails,
        $2p=\tfrac32>1$ works.
  \item \emph{Origin.} $f=0$ on both axes, so both partials are $0$. The
        remainder is $r(\mathbf{h})=|h_{1}h_{2}|\le\tfrac12\|\mathbf{h}\|^{2}$,
        so $|r|/\|\mathbf{h}\|\le\tfrac12\|\mathbf{h}\|\to0$: differentiable,
        with tangent plane $z=0$. \emph{At $(2,0)$.}
        $f(2,h)-f(2,0)=|2h|=2|h|$, and the quotient $2|h|/h=\pm2$ has no
        limit: $f_{y}(2,0)$ does not exist, so $f$ is not differentiable
        there. The graph is creased along both axes, except at the origin,
        where the two creases flatten out.
  \item \emph{Partials.} $f(h,0)=0$ and $f(0,h)=0$ for $h\neq0$, so both
        difference quotients are $0/h=0$: $f_{x}(0,0)=f_{y}(0,0)=0$.
        \emph{Lines.} On the $y$-axis $f=0$. On $y=mx$, $x\neq0$:
        \[
          f(x,mx)=\frac{mx^{3}}{x^{4}+m^{2}x^{2}}=\frac{mx}{x^{2}+m^{2}}
          \;\longrightarrow\;0\quad(x\to0),
        \]
        for $m\neq0$, and $f=0$ for $m=0$. Every straight line gives $0$.
        \emph{Parabola.} On $y=x^{2}$: $f=\dfrac{x^{4}}{x^{4}+x^{4}}=\tfrac12$
        for every $x\neq0$. So $f$ has no limit at the origin and is
        \textbf{not continuous} there, hence \textbf{not differentiable},
        although both partials exist and every line agrees.
  \item \emph{Partials.} $f=0$ on both axes, so $f_{x}(0,0)=f_{y}(0,0)=0$.
        \emph{Lines.} On $y=mx$:
        $f=\dfrac{2m^{2}x^{3}}{x^{2}+m^{4}x^{4}}=\dfrac{2m^{2}x}{1+m^{4}x^{2}}
        \to0$; on $x=0$, $f=0$. \emph{Parabola.} On $x=y^{2}$:
        $f=\dfrac{2y^{4}}{y^{4}+y^{4}}=1$ for $y\neq0$. Not continuous at the
        origin, hence not differentiable.
  \item (a) A smooth bowl, with no edges, creases or tips anywhere: it has a
        tangent plane at every point (horizontal at the bottom), so the
        function is differentiable everywhere. (b) A cone with a sharp tip:
        differentiable everywhere except at the tip, where every vertical
        cut through the tip is a V and no plane fits. (c) A pyramid whose
        four flat faces meet along two creases running over the axes:
        differentiable on the faces, not on the two creases (including
        the top point).
  \item The bowl $(x-1)^{2}+4y^{2}$ has no tip, so it is out. The level
        $c>0$ of $\sqrt{(x-1)^{2}+4y^{2}}$ is
        $\dfrac{(x-1)^{2}}{c^{2}}+\dfrac{y^{2}}{(c/2)^{2}}=1$, semi-axes $c$
        in $x$ and $c/2$ in $y$: twice as long in $x$, as in the picture.
        The third formula has semi-axes $c/2$ in $x$ and $c$ in $y$, the
        wrong way round. So (b) is $\sqrt{(x-1)^{2}+4y^{2}}$. At $(1,0)$:
        $f(1+h,0)-f(1,0)=\sqrt{h^{2}}=|h|$, and $|h|/h=\pm1$ has no limit;
        likewise $f(1,h)-f(1,0)=2|h|$. Neither partial exists, so $f$ is
        not differentiable at the tip, as the picture said.
\end{solutions}
